\chapter{Overview}

The first milestone is Hamilton's 1982 theorem: a closed $3$-manifold admitting
a metric of strictly positive Ricci curvature admits a metric of constant
positive sectional curvature. The terminal node is the spherical space form
theorem --- a closed $3$-manifold with finite fundamental group is
$S^3/\Gamma$ --- which Perelman proved in 2003 and which contains the
Poincar\'e conjecture as the case $\Gamma = 1$. Mathlib states that corollary,
unproved, in \texttt{Mathlib/Geometry/Manifold/PoincareConjecture.lean}.

The road is charted to the space form theorem rather than to geometrization on
purpose. Finite extinction time makes the long-time behaviour of the flow
irrelevant, which removes collapsing theory, graph manifolds, and the JSJ
decomposition from the dependency graph entirely --- see the end of Chapter
\ref{chap:spaceform}.

The obstacle is usually said to be analysis, and eventually it is --- but not
first. When this blueprint began, Riemannian curvature did not exist in Mathlib
at all: the library had manifolds, vector bundles, and covariant derivatives,
but no Riemann tensor, no $\Ric$, and not one occurrence of the string
\texttt{sectional}. A project planning around the parabolic gap therefore
stalls long before reaching it. Chapters \ref{chap:curvature} and
\ref{chap:levi-civita} close that first gap: $\Riem$, $\Ric$ and $\Sect$ are
now defined, and Hamilton's theorem is stated in Lean.

Since then the Levi-Civita connection has landed in Mathlib, and this
blueprint proves it is smooth (Theorem \ref{thm:levi-civita-smooth}), so
$\Ric$ is now a function of the metric and the flow is the textbook equation
$\partial_t \g = -2\Ric(\g)$ (Lemma \ref{lem:flow-textbook}).

The analysis gap itself is real and unmoved. Mathlib has no parabolic PDE
theory, no elliptic regularity, and no Sobolev spaces of bundle sections, so
short-time existence has nothing to stand on; Chapter \ref{chap:flow} marks
that region explicitly. It is also wider than usually described: the road needs
\emph{three} distinct parabolic theories --- Ricci flow, harmonic map flow, and
curve-shrinking flow --- and a second independent gap at Cheeger--Gromov
compactness (Chapter \ref{chap:kappa}).

Two branches are nevertheless closed end to end, both free of PDE: Chapter
\ref{chap:homog}, where the flow on left-invariant metrics is an ODE, and
Chapter \ref{chap:milnor}, Milnor's curvature formulas.

\section*{State of the art, September 2026}

Hamilton's theorem is no longer open. Chow, Liao and Qin
(\texttt{arXiv:2608.21502}, 21 August 2026) give a Lean formalization of it:
the endpoint \texttt{hamilton\_positive\_ricci} concludes
\emph{constant positive sectional curvature and spherical space form} from
\emph{closed connected three-manifold} and \emph{admits positive Ricci},
with no auxiliary hypotheses, no \texttt{sorry} and no project axiom. Their
library is roughly $1.9$ million lines. It closes, in the language of this
chart, Chapter \ref{chap:flow}'s short-time existence and maximal
continuation, Shi's estimates, both maximum principles,
Theorem \ref{thm:w-monotone}, Theorem \ref{thm:no-local-collapsing} and
Theorem \ref{thm:compactness} --- that is, \emph{both} of the analytic gaps
this overview names above, together with the whole of Chapter
\ref{chap:hamilton}.

So the frontier has moved, and it has moved past the two obstacles that were
supposed to be the hard ones. What is left between here and
Corollary \ref{cor:poincare} is no longer parabolic existence theory. It is,
in dependency order:

\begin{itemize}
  \item Theorem \ref{thm:hamilton-ivey}, the Hamilton--Ivey estimate --- an
    ODE invariant set plus a time-dependent tensor maximum principle, with no
    new analysis in it at all, and absent from their library;
  \item Definition \ref{def:reduced-volume} and
    Theorem \ref{thm:reduced-volume-monotone}, Perelman's
    $\mathcal{L}$-geometry --- two chapters of comparison geometry built on
    the flow, needing Jacobi fields and the second variation but no parabolic
    theory, and likewise absent;
  \item Theorem \ref{thm:kappa-classification} and
    Theorem \ref{thm:canonical-nbhd}, which consume the two above together
    with the compactness theory that now exists;
  \item Theorem \ref{thm:surgery-exists} and
    Theorem \ref{thm:finite-extinction}, which carry the two remaining
    parabolic prerequisites --- harmonic map flow and curve-shrinking flow.
\end{itemize}

The first two items are the ones reachable from what is proved here.

\chapter{Curvature}
\label{chap:curvature}

\begin{definition}[Curvature of a covariant derivative]
  \label{def:curvature}
  \lean{CovariantDerivative.curvature}
  \leanok
  For a covariant derivative $\nab$ on the tangent bundle of $\M$,
  \[ \Riem(X,Y)Z = \nab_X \nab_Y Z - \nab_Y \nab_X Z - \nab_{[X,Y]} Z. \]
\end{definition}

\begin{lemma}[Antisymmetry]
  \label{lem:curvature-antisymm}
  \lean{CovariantDerivative.curvature_antisymm}
  \leanok
  \uses{def:curvature}
  $\Riem(X,Y)Z = -\Riem(Y,X)Z$.
\end{lemma}

\begin{proof}
  \leanok
  Immediate from $[Y,X] = -[X,Y]$ and linearity of $\nab_{(-)} Z$ in the
  direction slot.
\end{proof}

\begin{lemma}[First Bianchi identity]
  \label{lem:bianchi-first}
  \lean{CovariantDerivative.bianchi_first}
  \leanok
  \uses{def:curvature, thm:levi-civita}
  For a torsion-free connection,
  $\Riem(X,Y)Z + \Riem(Y,Z)X + \Riem(Z,X)Y = 0$.
\end{lemma}

\begin{proof}
  \leanok
  \textbf{Proved}, for a $C^1$ connection and $C^2$ vector fields, with every
  differentiability side condition discharged.

  Regroup the cyclic sum
  using torsion-freeness on each pair,
  \[ \textstyle\sum \Riem(X,Y)Z = \nab_X [Y,Z] - \nab_{[Y,Z]} X
     + \nab_Y [Z,X] - \nab_{[Z,X]} Y + \nab_Z [X,Y] - \nab_{[X,Y]} Z, \]
  apply torsion-freeness once more to each difference to obtain
  $[X,[Y,Z]] + [Y,[Z,X]] + [Z,[X,Y]]$, and close with Jacobi
  (\texttt{VectorField.leibniz\_identity\_mlieBracket\_apply}).

  Both ingredients are in Mathlib as of August 2026. The $C^2$ hypotheses are
  essential rather than cosmetic: \texttt{torsion\_eq\_zero\_iff} gives
  $\nab_X Y - \nab_Y X = [X,Y]$ only for \texttt{MDiffAt} fields, and the
  second application is to the pair $(X, [Y,Z])$, so the bracket itself must be
  differentiable.

  (An earlier version of this blueprint claimed the identity was blocked on
  $Z$-slot tensoriality and \texttt{mkHom}$_3$. Those are prerequisites for
  packaging curvature as a \emph{tensor}, not for the operator identity stated
  here.)
\end{proof}

\begin{lemma}[Second Bianchi identity]
  \label{lem:bianchi-second}
  \lean{CovariantDerivative.bianchi_second}
  \leanok
  \uses{def:curvature, lem:bianchi-first}
  For a torsion-free connection,
  $(\nab_X \Riem)(Y,Z)W + (\nab_Y \Riem)(Z,X)W + (\nab_Z \Riem)(X,Y)W = 0$,
  where $(\nab_X \Riem)(Y,Z)W = \nab_X(\Riem(Y,Z)W) - \Riem(\nab_X Y, Z)W
  - \Riem(Y, \nab_X Z)W - \Riem(Y,Z)\nab_X W$
  (\texttt{CovariantDerivative.covCurvature}).
\end{lemma}

\begin{proof}
  \leanok
  \textbf{Proved}, for a $C^2$ connection, $C^2$ fields $X, Y, Z$ and a $C^3$
  field $W$, with every differentiability side condition discharged. No metric
  is involved.

  Expand everything in iterated covariant derivatives of $W$. The twelve
  third-order terms $\nab\nab\nab W$ cancel in pairs, as do the six terms
  $\nab_X \nab_{[Y,Z]} W$ and $\nab_{[Y,Z]} \nab_X W$ once torsion-freeness has
  turned $\nab_{\nab_X Y} - \nab_{\nab_Y X}$ into $\nab_{[X,Y]}$ in the
  direction slot and $\nab_Z \nab_{\nab_X Y} W - \nab_Z \nab_{\nab_Y X} W$
  into $\nab_Z \nab_{[X,Y]} W$ as sections. What survives is
  $\nab_W$ applied to $[[X,Y],Z] + [[Y,Z],X] + [[Z,X],Y]$, which is Jacobi.

  The regularity is what makes the sections differentiable: a $C^2$
  connection sends the $C^3$ section $W$ to a $C^2$ section $\nab_Z W$
  (\texttt{contMDiff\_cov\_apply}), so $\nab_Y \nab_Z W$ is differentiable and
  $\nab_X$ of it is meaningful. As for the first identity, brackets of $C^2$
  fields are differentiable, which the second application of torsion-freeness
  needs.
\end{proof}

\begin{lemma}[The Lie bracket is a derivation]
  \label{lem:derivation}
  \lean{VectorField.mlieBracket_apply_fun}
  \leanok
  $[V,W]f = V(Wf) - W(Vf)$, for $f$ of class $C^2$ at the point and $V,W$
  differentiable there, on any manifold --- boundary and corners allowed.
\end{lemma}

\begin{proof}
  \leanok
  Absent from Mathlib in both the vector-space and the manifold setting. Over a
  vector space the two second-derivative terms cancel by symmetry of the second
  derivative, leaving exactly the bracket. Transporting to a manifold turns on one
  observation: $\mathrm{d}f$ at a nearby point $y$ is defined through the chart
  \emph{at $y$}, so one first re-expresses it through the single chart at $x$,
  uniformly over that chart's source. The outer derivative is then an ordinary
  \texttt{fderivWithin} and the vector-space identity transports verbatim, with
  $\mathrm{range}\,I$ supplying the side conditions the \texttt{Within} calculus
  needs. No manifold-level second-derivative symmetry is required.
\end{proof}

\begin{lemma}[Tensoriality of the curvature operator]
  \label{lem:tensoriality}
  \lean{CovariantDerivative.curvature_smul_left}
  \leanok
  \uses{def:curvature, lem:derivation}
  $\Riem$ is $C^\infty(\M)$-linear in each slot.
\end{lemma}

\begin{proof}
  \leanok
  First slot: expanding $\nab_{fX}Z = f\,\nab_X Z$ leaves a $(Yf)\,\nab_X Z$ term
  from Leibniz, and $[fX,Y] = f[X,Y] - (Yf)X$ leaves a matching one from the
  bracket, with opposite sign; they cancel. The second slot follows by
  antisymmetry. The third leaves a residue $[\,X(Yf) - Y(Xf) - [X,Y]f\,]\,Z$,
  which vanishes by Lemma~\ref{lem:derivation}.
\end{proof}

\chapter{The Levi-Civita connection}
\label{chap:levi-civita}

Mathlib has the Levi-Civita connection as of September 2026: mathlib4
PR \#36845 (Massot, Rothgang, Macbeth) added
\texttt{CovariantDerivative/LeviCivita.lean} with the Koszul formula,
uniqueness, a construction by the musical isomorphism, and the proof that the
construction is compatible and torsion-free. This chapter records what that
gives the road, and the one thing it does not give yet.

\begin{theorem}[Existence]
  \label{thm:levi-civita}
  \lean{RicciFlowBlueprint.exists_leviCivita}
  \leanok
  For a $C^1$ Riemannian metric $\g$ on a $C^2$ manifold there is a torsion-free
  covariant derivative on $TM$ compatible with $\g$.
\end{theorem}

\begin{proof}
  \leanok
  Mathlib's \texttt{leviCivitaConnection I M}: the Koszul expression
  $(X, Z) \mapsto \langle \nab_X Y, Z\rangle$ is tensorial in $X$ and $Z$ for
  differentiable $Y$, so it is a $(2,0)$-tensor at each point; the musical
  isomorphism turns it into an endomorphism $\nab Y$, and that is the
  connection. Compatibility and torsion-freeness are then computations with the
  Koszul formula. No local frames are involved.
\end{proof}

\begin{theorem}[Uniqueness]
  \label{thm:levi-civita-unique}
  \lean{RicciFlowBlueprint.leviCivita_unique}
  \leanok
  Two torsion-free metric-compatible covariant derivatives agree on every vector
  field that is differentiable at the point.
\end{theorem}

\begin{proof}
  \leanok
  The Koszul formula determines $\langle \nab_X Y, Z\rangle$ for differentiable
  $X, Y, Z$ without reference to $\nab$, and a vector is determined by its inner
  products against differentiable sections.

  This --- not ``there is exactly one'' --- is the fundamental theorem. A
  \texttt{CovariantDerivative} is unconstrained on sections that are not
  differentiable at the point: the Leibniz and additivity laws, compatibility
  and the torsion all quantify over differentiable sections only, so a
  Levi-Civita connection can be altered on non-differentiable sections and stay
  Levi-Civita. The $\exists!$ statement this chapter carried until now, as its
  only \texttt{sorry}, was therefore not provable as stated. Nothing downstream
  needed it: every curvature computation evaluates $\nab$ on differentiable
  sections.
\end{proof}

\begin{lemma}[Curvature and Ricci are pointwise in the last slot]
  \label{lem:curvature-pointwise}
  \lean{CovariantDerivative.curvature_congr_third, CovariantDerivative.curvature_smul_third, CovariantDerivative.curvature_smul_third', CovariantDerivative.ricci_congr_snd_of_eq, CovariantDerivative.ricciAt, CovariantDerivative.inner_curvature_eq_zero_of_dep, RicciFlowBlueprint.hasPositiveRicciLC_iff_tangent, RicciFlowBlueprint.hasConstSecLC_iff_mul}
  \leanok
  \uses{lem:global-extension, def:curvature, def:ricci}
  $\Riem(X,Y)Z$ at $x$ depends only on $Z(x)$, and hence $\Ric(X,Y)(x)$ depends
  only on $X(x)$ and $Y(x)$. Consequently Ricci curvature is a genuine function
  of tangent vectors, $\texttt{ricciAt}\ x : T_x\M \to T_x\M \to \R$, and
  positivity of Ricci can be stated pointwise:
  \[
    \forall x,\ \forall v \in T_x\M,\ v \ne 0 \implies 0 < \Ric_x(v,v).
  \]

  \textbf{Proved.} Expand $Z$ in a local frame, valid near $x$
  (\texttt{OrthonormalFrame.lean}); globalise the frame fields and the
  coefficient functions by Lemma \ref{lem:global-extension}; apply
  $C^\infty(\M)$-linearity in the third slot
  (\texttt{curvature\_sum\_smul\_third}, whose twelve side conditions are
  discharged from globally $C^2$ data in \texttt{curvature\_smul\_third}); the
  coefficients at $x$ are determined by $Z(x)$. Locality in the third slot
  (Lemma \ref{lem:global-extension}) is what lets the frame expansion, which
  holds only near $x$, be used at all.

  Only \emph{pointwise} regularity of the first two slots is needed, which is
  what makes the result applicable to \texttt{FiberBundle.extend} and hence to
  the trace defining $\Ric$.

  The same machinery gives the well-definedness of the sectional curvature on a
  general manifold (\texttt{sectionalCurvature\_congr'}, whose hypothesis
  \texttt{h3} was previously assumed), and, with the vanishing of
  $\langle \Riem(X,Y)Y, X\rangle$ on a linearly dependent pair
  (\texttt{inner\_curvature\_eq\_zero\_of\_dep}, by the equality case of
  Cauchy--Schwarz and antisymmetry in the first two slots), the
  \emph{multiplied-out} form of constant sectional curvature
  (\texttt{hasConstSecLC\_iff\_mul}):
  \[
    \langle \Riem(X,Y)Y, X\rangle
      = k\,\bigl(|X|^2|Y|^2 - \langle X,Y\rangle^2\bigr),
  \]
  with no division and no nondegeneracy hypothesis, which is how Chow--Liao--Qin
  state it.

  This is the form Chow--Liao--Qin's \texttt{positiveRicciMetric} takes
  (arXiv 2608.21502), where it comes from a tensor-bundle construction; here it
  is derived instead from the frame and a bump function. It closes the last
  statement-level divergence recorded in
  \texttt{notes/hamilton-statement-comparison.md}. Note the result is a
  congruence, not \texttt{TensorialAt}: that structure demands its identities
  for merely differentiable sections, which the third slot cannot satisfy.
\end{lemma}

\begin{lemma}[Global extension of a tangent vector]
  \label{lem:global-extension}
  \lean{RicciFlowBlueprint.exists_contMDiff_eventuallyEq, RicciFlowBlueprint.exists_contMDiff_extension, RicciFlowBlueprint.forall_contMDiff_iff_forall_tangent, CovariantDerivative.curvature_congr_third_of_eventuallyEq}
  \leanok
  On a Hausdorff manifold, every tangent vector $v \in T_x\M$ is the value at $x$
  of a globally $C^k$ vector field. Consequently, quantifying a property of
  $X(x)$ over globally $C^2$ fields $X$ is the same as quantifying it over
  tangent vectors.

  \textbf{Proved.} Mathlib's \texttt{FiberBundle.extend} produces a section
  through $v$ that is $C^k$ only on a neighbourhood of $x$
  (\texttt{exists\_contMDiffOn\_extend}). Multiplying by a smooth bump function
  whose closed support lies inside that neighbourhood
  (\texttt{SmoothBumpFunction.nhds\_basis\_support}) gives a global field, still
  equal to $v$ at $x$ because the bump is $1$ at its centre. Hausdorffness is
  needed for the bump function to be smooth and is exactly the hypothesis whose
  absence made the curvature predicates vacuously satisfiable; see
  \texttt{notes/hamilton-statement-comparison.md}.

  The general form is \texttt{exists\_contMDiff\_eventuallyEq}: a section that is
  $C^k$ on a neighbourhood of $x$ agrees near $x$ with a globally $C^k$ section,
  with \texttt{exists\_contMDiff\_eventuallyEq\_fun} the scalar companion for
  functions.
  This closes the long-standing gap "no global $C^2$ extension of a tangent
  vector".

  Alongside it, \texttt{curvature\_congr\_third\_of\_eventuallyEq}: $\Riem(X,Y)Z$
  at $x$ depends only on $Z$ \emph{near} $x$, since $\nab$ is local on
  differentiable sections.

  These two are what the pointwise statement below is built from.
\end{lemma}

\begin{theorem}[Curvature is well defined]
  \label{thm:ricci-well-defined}
  \lean{RicciFlowBlueprint.ricci_eq_of_isLeviCivita}
  \leanok
  \uses{thm:levi-civita-unique, def:ricci, def:sectional}
  Any two $C^1$ Levi-Civita connections have the same Riemann tensor on fields
  $X, Y$ differentiable at the point and $C^2$ field $Z$, hence the same Ricci
  tensor and the same sectional curvature on $C^2$ fields.
\end{theorem}

\begin{proof}
  \leanok
  $R(X,Y)Z = \nab_X \nab_Y Z - \nab_Y \nab_X Z - \nab_{[X,Y]} Z$. The inner
  derivatives agree as sections by uniqueness, since $Z$ is differentiable
  everywhere; the outer derivatives then agree by uniqueness again, since
  $\nab_Y Z$ is differentiable at $x$ when $\nab$ is $C^1$ and $Z$ is $C^2$.
  For $\Ric$, the trace is over the endomorphism produced by
  \texttt{TensorialAt.mkHom}, whose value on a vector is $R$ evaluated on a
  differentiable extension of it. Sectional curvature is a quotient of inner
  products of $R$.

  This is what makes the existential quantification over connections in
  Chapters \ref{chap:flow} and \ref{chap:hamilton} well-posed: the
  predicates there are statements about the metric, not about a choice of
  connection.
\end{proof}

\begin{theorem}[Smoothness of the Levi-Civita connection]
  \label{thm:levi-civita-smooth}
  \lean{RicciFlowBlueprint.contMDiffCovariantDerivative_leviCivitaConnection}
  \leanok
  \uses{thm:levi-civita}
  On a $C^\omega$ manifold with a $C^{k+1}$ Riemannian metric, Mathlib's
  \texttt{leviCivitaConnection} is a $C^k$ covariant derivative: for every
  $C^{k+1}$ vector field $Y$, the section $x \mapsto \nab Y(x)$ of
  $\operatorname{End}(TM)$ is $C^k$.
\end{theorem}

\begin{proof}
  \leanok
  This is the result Mathlib's \texttt{LeviCivita.lean} defers to ``future
  PRs''. Work in a trivialisation $e$ of $TM$ at $x_0$ with the local frame
  $s_i = e^{-1}(\cdot, b_i)$ for a basis $b$ of the model space. Three
  criteria, each proved separately:
  \begin{enumerate}
  \item A map into $\operatorname{Hom}(E, F)$ is $C^k$ iff its values on a
    basis of $E$ are (\texttt{contMDiffAt\_clm\_of\_basis}).
  \item A section $w$ of a Riemannian bundle with $C^k$ metric is $C^k$ at
    $x_0$ if the scalar functions $\langle w, s_j\rangle$ are
    (\texttt{contMDiffAt\_section\_of\_inner\_localFrame}): the coordinates
    of $w$ are $G^{-1}(\langle w, s_j\rangle)_j$ for the Gram operator $G$ of
    the frame, $G$ is smooth and invertible on the base set, and
    \texttt{ContinuousLinearMap.inverse} is smooth at invertible points.
  \item The Koszul expression for $\langle \nab_{s_i} Y, s_j\rangle$ is
    $C^k$: derivatives of $C^{k+1}$ inner products along $C^{k}$ fields
    (\texttt{contMDiffAt\_mvfderiv\_apply}) and inner products of $C^{k+1}$
    fields with brackets of $C^{k+1}$ fields, which are $C^k$.
  \end{enumerate}
  Then $\nab_{s_i} Y$ is $C^k$ by (2) and (3), and $\nab Y$ in coordinates
  is determined by the $\nab_{s_i} Y$ by (1). No orthonormal frames are needed;
  the Gram inversion replaces them.
\end{proof}

\begin{definition}[Curvature of a metric]
  \label{def:ricci-of-metric}
  \lean{RicciFlowBlueprint.ricciOfMetric}
  \leanok
  \uses{thm:levi-civita-smooth, def:ricci, def:sectional}
  For a $C^2$ metric $\g$, $\Ric(\g)$ is the Ricci curvature of its
  Levi-Civita connection, and likewise $\Sect(\g)$. By
  Theorem \ref{thm:levi-civita-smooth} the connection is $C^1$, so these are
  genuine tensors and not the junk value that \texttt{ricci} returns on a
  connection not known to be $C^1$. This is the function
  $\g \mapsto \Ric(\g)$ that DeTurck's trick and the curvature evolution
  equations need.
\end{definition}

\begin{definition}[Ricci and scalar curvature]
  \label{def:ricci}
  \lean{CovariantDerivative.ricci}
  \uses{def:curvature, thm:levi-civita, lem:tensoriality}
  $\Ric(X,Y) = \operatorname{tr}(Z \mapsto \Riem(Z,X)Y)$ and
  $\scal = \operatorname{tr}_\g \Ric$.

  \leanok
  \textbf{Defined.} The trace needs $\Riem$ to descend from vector \emph{fields}
  to vectors, and the criterion is $C^\infty(\M)$-linearity in each slot ---
  Lemma~\ref{lem:tensoriality}, together with additivity
  (\texttt{curvature\_add\_left/middle/right}). Since the trace is taken in the
  \emph{first} slot only, Mathlib's one-slot \texttt{TensorialAt.mkHom}
  suffices and the \texttt{mkHom}$_3$ this blueprint previously called for is
  not needed --- and would not have worked: its hypotheses quantify over merely
  \texttt{MDiffAt} sections, whereas first-slot tensoriality of $\Riem$ needs
  the third-slot field to be $C^2$.
\end{definition}

\begin{lemma}[Trace of the first Bianchi identity]
  \label{lem:trace-bianchi}
  \lean{CovariantDerivative.ricci_sub_ricci_swap}
  \leanok
  \uses{def:ricci, lem:bianchi-first}
  For a torsion-free connection,
  $\Ric(X,Y) - \Ric(Y,X) = -\operatorname{tr} \Riem(X,Y)$.
\end{lemma}

\begin{proof}
  \leanok
  Trace the cyclic sum $\Riem(X,Y)Z + \Riem(Y,Z)X + \Riem(Z,X)Y = 0$ in $Z$:
  the first term contributes $\operatorname{tr}\Riem(X,Y)$, the second
  $-\Ric(Y,X)$ by antisymmetry, the third $\Ric(X,Y)$ by definition.

  \textbf{Note.} $\Ric$ is \emph{not} symmetric for a general torsion-free
  connection; symmetry is the corollary
  (\texttt{ricci\_symm\_of\_traceless}) obtained when
  $\operatorname{tr}\Riem(X,Y) = 0$, which holds for a metric connection
  because $\Riem(X,Y)$ is then skew-adjoint. An earlier plan for this blueprint
  asserted symmetry from torsion-freeness alone. That is false.
\end{proof}

\begin{definition}[The Ricci form]
  \label{def:ricci-form}
  \lean{CovariantDerivative.ricciForm, CovariantDerivative.ricciForm_apply_field,
    CovariantDerivative.ricciAt_add_left, CovariantDerivative.ricciAt_smul_left,
    CovariantDerivative.ricciAt_add_right, CovariantDerivative.ricciAt_smul_right}
  \leanok
  \uses{def:ricci, lem:curvature-pointwise, lem:global-extension}
  $\Ric$ as an honest continuous bilinear form $\Ric_x : T_x\M \times T_x\M \to
  \R$ on a general manifold.

  \textbf{Proved} (\texttt{RicciForm.lean}). Lemma
  \ref{lem:curvature-pointwise} makes $\Ric(X,Y)(x)$ depend on $X_x$ and $Y_x$
  alone, so $\Ric_x$ is a well-defined \emph{function} of two tangent vectors
  (\texttt{ricciAt}); each of its four linearity laws is the corresponding
  field-level law (\texttt{ricci\_add\_left} and friends) applied to global
  $C^2$ extensions, which exist by Lemma \ref{lem:global-extension}.
  Continuity is automatic in finite dimension. Two years' worth of the
  ``$\Ric$ is not a tensor here'' caveats in this blueprint dissolve at this
  node.
\end{definition}

\begin{lemma}[Symmetries of the Riemann tensor]
  \label{lem:curvature-symmetries}
  \lean{CovariantDerivative.inner_curvature_pair_symm,
    CovariantDerivative.inner_bianchi_first,
    CovariantDerivative.inner_curvature_left_skew,
    CovariantDerivative.inner_curvature_right_skew'}
  \leanok
  \uses{def:curvature, lem:bianchi-first, thm:levi-civita}
  Write $\Rm(X,Y,Z,W) = \langle \Riem(X,Y)Z, W\rangle$ for a metric torsion-free
  connection. Then
  \[
    \Rm(X,Y,Z,W) = -\Rm(Y,X,Z,W), \qquad
    \Rm(X,Y,Z,W) = -\Rm(X,Y,W,Z),
  \]
  \[
    \Rm(X,Y,Z,W) + \Rm(Y,Z,X,W) + \Rm(Z,X,Y,W) = 0, \qquad
    \Rm(X,Y,Z,W) = \Rm(Z,W,X,Y).
  \]

  \textbf{Proved} (\texttt{CurvatureSymm.lean}). The first three were already
  here: antisymmetry in the first pair needs no hypothesis at all
  (\texttt{curvature\_antisymm}), antisymmetry in the second pair is metric
  compatibility (\texttt{inner\_curvature\_right\_skew}), and the cyclic
  identity is Lemma \ref{lem:bianchi-first} paired against a vector --- note the
  fourth slot needs no regularity, the identity being a vector identity.

  \textbf{Pair symmetry} is the new one, and it is what makes $\Riem$ a
  symmetric operator on $\Lambda^2 T_x\M$ --- the form Hamilton's pinching
  estimates and Uhlenbeck's trick are stated in --- and what the second
  contraction of the second Bianchi identity will need.
\end{lemma}

\begin{proof}
  \leanok
  The classical octahedron argument. In the six unknowns
  \[
    p = \Rm(X,Y,Z,W),\ q = \Rm(Z,W,X,Y),\ r = \Rm(Y,Z,W,X),
  \]
  \[
    s = \Rm(X,Z,W,Y),\ t = \Rm(Y,W,X,Z),\ u = \Rm(W,X,Y,Z),
  \]
  the four cyclic identities on $(X,Y,Z;W)$, $(Y,Z,W;X)$, $(Z,W,X;Y)$ and
  $(W,X,Y;Z)$, after every term has been brought to one of these six by the two
  antisymmetries, read
  \[
    p - r + s = 0,\quad r - q + t = 0,\quad q - u + s = 0,\quad u - p + t = 0 .
  \]
  Adding the first two and the last two gives $p - q + s + t = 0$ and
  $q - p + s + t = 0$; subtracting, $2(p-q) = 0$. In Lean the six links are
  supplied as hypotheses and \texttt{linarith} does the elimination.
\end{proof}

\begin{lemma}[Constant sectional curvature determines $\Rm$]
  \label{lem:const-sec-determines-rm}
  \lean{CovariantDerivative.inner_curvature_eq_of_const_sec,
    CovariantDerivative.inner_curvature_eq_of_const_sec_norm,
    CovariantDerivative.curvatureDefect,
    CovariantDerivative.ricci_eq_of_const_sec,
    CovariantDerivative.scalarCurvatureAt_eq_of_const_sec,
    RicciFlowBlueprint.hasConstSecLC_tensor,
    RicciFlowBlueprint.hasConstSecLC_ricci,
    RicciFlowBlueprint.hasConstSecLC_scalar}
  \leanok
  \uses{def:sectional, lem:curvature-symmetries, lem:bianchi-first, def:ricci-form, def:scalar}
  If every sectional curvature at $x$ equals $k$, then for \emph{all} four
  arguments
  \[
    \langle \Riem(X,Y)Z, W\rangle
      = k\big(\langle Y,Z\rangle\langle X,W\rangle
             - \langle X,Z\rangle\langle Y,W\rangle\big).
  \]

  \textbf{Proved} (\texttt{ConstantCurvature.lean}). This is a strengthening,
  not a parity fix. Chow--Liao--Qin's
  \texttt{admitsConstantPositiveSectionalCurvature} is the multiplied-out
  \emph{sectional} identity $\Rm(X,Y,Y,X) = c(\g(X,X)\g(Y,Y) - \g(X,Y)^2)$,
  which \texttt{hasConstSecLC\_iff\_mul} already matches; the divergence
  recorded in \texttt{notes/hamilton-statement-comparison.md} was closed before
  this lemma existed. What the lemma adds is that the sectional identity pins
  down the \emph{whole} $(0,4)$ tensor, so everything downstream of constant
  curvature --- the spherical space form conclusion above all --- can use $\Rm$
  directly instead of re-deriving it.

  \textbf{Corollaries}, by tracing once and twice
  (\texttt{ricci\_eq\_of\_const\_sec},
  \texttt{scalarCurvatureAt\_eq\_of\_const\_sec}):
  \[
    \Ric = (n-1)k\,\g, \qquad \scal = n(n-1)k .
  \]
  Both are stated over an arbitrary orthonormal basis of $T_x\M$, with $n$ its
  cardinality, so no identification of $\operatorname{finrank}$ with the
  dimension of the model is needed. Tracing uses Parseval
  (\texttt{OrthonormalBasis.sum\_inner\_mul\_inner}) exactly once, on
  $\sum_i \langle X, e_i\rangle \langle e_i, Y\rangle = \langle X, Y\rangle$;
  the first slot of $\Ric$ is fed the global $C^2$ extensions of the basis
  vectors, which is what \texttt{extendTwo} is for.

  \textbf{Applied to the predicate.} \texttt{hasConstSecLC\_tensor},
  \texttt{hasConstSecLC\_ricci} and \texttt{hasConstSecLC\_scalar} instantiate
  all three at \texttt{HasConstSecLC I M k} and the canonical Levi-Civita
  connection, via \texttt{hasConstSecLC\_iff\_mul}. So the conclusion of
  Theorem \ref{thm:hamilton} delivers the full $(0,4)$ tensor, and the metric it
  produces is Einstein with $\scal = n(n-1)k$ --- which is what the spherical
  space form direction will want to consume.
\end{lemma}

\begin{proof}
  \leanok
  Polarisation, twice. Let
  \[
    T(A,B,C,D) = \langle \Riem(A,B)C, D\rangle
      - k\big(\langle B,C\rangle\langle A,D\rangle
             - \langle A,C\rangle\langle B,D\rangle\big)
  \]
  (\texttt{curvatureDefect}). The model term has every symmetry the curvature
  has, so by Lemma \ref{lem:curvature-symmetries} $T$ is antisymmetric in each
  pair, symmetric under exchanging the pairs, and satisfies the first Bianchi
  identity; and the hypothesis says exactly $T(A,B,B,A) = 0$.

  Expanding $T(A+C, B, B, A+C) = 0$ and using $T(C,B,B,A) = T(A,B,B,C)$ --- pair
  symmetry followed by the two antisymmetries --- gives $2T(A,B,B,C) = 0$, so
  $T(A,B,B,C) = 0$ for all $A,B,C$. Expanding $T(A, B+C, B+C, D) = 0$ then gives
  $T(A,B,C,D) + T(A,C,B,D) = 0$: $T$ is antisymmetric in its middle two slots as
  well. A tensor antisymmetric in three consecutive slots has all three terms of
  the first Bianchi identity equal, so $3T(X,Y,Z,W) = 0$.
\end{proof}

\begin{theorem}[Ricci is symmetric]
  \label{thm:ricci-symm}
  \lean{CovariantDerivative.ricci_symm, CovariantDerivative.ricciAt_symm,
    CovariantDerivative.ricciForm_symm, CovariantDerivative.curvatureEndoAt,
    CovariantDerivative.trace_curvatureEndoAt_eq_zero}
  \leanok
  \uses{lem:trace-bianchi, def:ricci-form, thm:levi-civita, lem:curvature-symmetries}
  For a metric torsion-free connection --- in particular for the Levi-Civita
  connection --- $\Ric(X,Y) = \Ric(Y,X)$ at every point of a general manifold,
  with no further hypothesis.

  \textbf{Torsion-freeness alone is not enough.} Lemma \ref{lem:trace-bianchi}
  says the antisymmetric part of $\Ric$ is exactly $-\operatorname{tr}(v \mapsto
  \Riem(X,Y)v)$, and that trace vanishes because the connection is \emph{metric}.
  Symmetry is a corollary of skew-adjointness, not of the first Bianchi identity.
\end{theorem}

\begin{proof}
  \leanok
  Lemma \ref{lem:trace-bianchi} carried two hypotheses it could not then
  discharge: that global $C^2$ fields realise every tangent vector, and that the
  endomorphism $v \mapsto \Riem(X,Y)v$ exists at all. Both are now theorems ---
  Lemma \ref{lem:global-extension} and Lemma \ref{lem:curvature-pointwise} ---
  so \texttt{curvatureEndoAt} is an unconditional linear endomorphism of
  $T_x\M$, its linearity being tensoriality of the curvature's third slot on
  global extensions. Its trace over an orthonormal basis is $\sum_i \langle e_i,
  \Riem(X,Y)e_i \rangle$, and each term vanishes because $\Riem(X,Y)$ is
  skew-adjoint for a metric connection (\texttt{inner\_curvature\_self\_right}).
  Transporting through Definition \ref{def:ricci-form} gives symmetry of the
  Ricci form itself.
\end{proof}

\begin{definition}[Scalar curvature]
  \label{def:scalar}
  \lean{CovariantDerivative.scalarCurvatureAt,
    CovariantDerivative.scalarCurvatureAt_eq_sum_basis,
    CovariantDerivative.scalarCurvatureWith_congr',
    CovariantDerivative.scalarCurvature}
  \leanok
  \uses{def:ricci, def:ricci-form}
  $\scal = \operatorname{tr}_\g \Ric$, the metric trace of the Ricci tensor:
  $\scal(x) = \sum_i \Ric(e_i, e_i)$ over an orthonormal basis of $T_x\M$.

  \textbf{Now a genuine trace on a general manifold}
  (\texttt{scalarCurvatureAt}), and frame-independence is \emph{free}:
  $\sum_i B(e_i,e_i)$ is basis-independent for any bilinear $B$, and by
  Definition \ref{def:ricci-form} $\Ric_x$ is one
  (\texttt{scalarCurvatureAt\_eq\_sum\_basis}, no hypotheses).

  This supersedes the earlier architecture, which is kept for the
  correspondence: the frame-relative \texttt{scalarCurvatureWith}, whose
  frame-independence \texttt{scalarCurvatureWith\_congr} had to assume
  pointwise dependence of the curvature's third slot in a form
  (\texttt{MDiffAt} in the first slot) that Lemma
  \ref{lem:curvature-pointwise} does not supply, and the hypothesis-free
  \texttt{scalarCurvature} on the model space. Going through the bilinear form
  discharges the hypothesis on any manifold
  (\texttt{scalarCurvatureWith\_congr'}).
\end{definition}

\begin{lemma}[Ricci as an orthonormal-frame sum]
  \label{lem:ricci-frame-sum}
  \lean{CovariantDerivative.ricci_eq_sum_inner_curvature}
  \leanok
  \uses{def:ricci}
  $\Ric(X,Y)(x) = \sum_i \langle b_i, \Riem(b_i, X)Y(x)\rangle$ for any
  orthonormal basis of the fibre --- hypothesis-free, on a general manifold.
  Raising an index is never needed: the frame sum with Parseval replaces the
  musical isomorphisms entirely.
\end{lemma}

\begin{definition}[Sectional curvature]
  \label{def:sectional}
  \lean{CovariantDerivative.sectionalCurvature}
  \leanok
  \uses{def:curvature, lem:tensoriality, thm:levi-civita}
  For a $2$-plane spanned by tangent vectors $X, Y$ at a point,
  \[ \Sect(X,Y) = \frac{\langle \Riem(X,Y)Y, X\rangle}
       {|X|^2|Y|^2 - \langle X,Y\rangle^2}, \]
  and the metric has \emph{constant} sectional curvature $k$ if $\Sect \equiv k$
  on every $2$-plane at every point.

  \textbf{Defined}, together with its invariance under change of basis of the
  spanned $2$-plane (\texttt{sectionalCurvature\_basis\_change}), which is the
  actual content. Absent from Mathlib entirely --- the string \texttt{sectional}
  does not occur anywhere in the library. This is the \emph{conclusion} side of
  Hamilton's theorem, and it is undefined for the same reason the hypothesis
  was: well-definedness on $2$-planes rather than on pairs of vector fields is
  exactly Lemma~\ref{lem:tensoriality}. Unlike $\Ric$ it needs no trace, so it
  requires the pointwise tensor but not \texttt{mkHom}$_3$.
\end{definition}

\chapter{Ricci flow on left-invariant metrics}
\label{chap:homog}

\textbf{This branch is closed end to end.} It is the one part of the road with
no PDE in it, and everything below is proved: no \texttt{sorry}, no
\texttt{proof\_wanted}, all axiom-free.

A left-invariant metric on a Lie group $G$ is determined by an inner product on
$\mathfrak{g}$, so $\partial_t \g = -2\Ric(\g)$ is an ODE on a
finite-dimensional space and Mathlib's Picard--Lindel\"of suffices. The
development lives purely at the level of a finite-dimensional real normed space
carrying a bracket, and shares \emph{no} machinery with the rest of this
blueprint --- no \texttt{CovariantDerivative}, no vector bundles, no
\texttt{RiemannianBundle}. Neither the analysis gap nor
Theorem~\ref{thm:levi-civita} is on its path.

Two representation choices earned this. The bracket is carried as an explicit
map rather than through \texttt{[LieRing]}: that class and
\texttt{NormedAddCommGroup} both extend \texttt{AddCommGroup}, an unresolved
diamond with no precedent in Mathlib, so with a norm in play the bracket must
leave the instance system. A transport bridge restates the results for genuine
Lie algebras, so no generality is lost. And the metric is carried as a map into
the dual, which makes smoothness a chain of
\texttt{contDiffAt\_map\_inverse} and composition lemmas rather than
coordinates and Cramer's rule.

\begin{definition}[The Koszul connection]
  \label{def:koszul-li}
  \lean{RicciFlowBlueprint.LeftInvariant.koszul}
  \leanok
  For left-invariant fields the Koszul formula
  \[ 2\langle \nab_X Y, Z\rangle = \langle [X,Y],Z\rangle - \langle [Y,Z],X\rangle
       + \langle [Z,X],Y\rangle \]
  determines $\nab$ algebraically, so it can be \emph{defined} rather than shown
  to exist.
\end{definition}

\begin{lemma}[Torsion-free and metric-compatible]
  \label{lem:koszul-li-props}
  \lean{RicciFlowBlueprint.LeftInvariant.koszul_sub_koszul}
  \leanok
  \uses{def:koszul-li}
  $\nab_X Y - \nab_Y X = [X,Y]$, and
  $\langle \nab_X Y, Z\rangle + \langle \nab_X Z, Y\rangle = 0$.
\end{lemma}

\begin{theorem}[Levi-Civita uniqueness, left-invariant case]
  \label{thm:levi-civita-li}
  \lean{RicciFlowBlueprint.LeftInvariant.eq_koszul_of_torsionFree_of_compat}
  \leanok
  \uses{lem:koszul-li-props}
  Any torsion-free, metric-compatible bilinear connection equals the Koszul
  connection.

  This is the fundamental theorem of Riemannian geometry in the left-invariant
  world --- the statement Theorem~\ref{thm:levi-civita} still leaves open at
  manifold level, discharged here because the algebra needs no existence
  theory.
\end{theorem}

\begin{lemma}[The Ricci field is smooth]
  \label{lem:ricci-lie-ode}
  \lean{RicciFlowBlueprint.LeftInvariant.contDiffAt_ricciField}
  \leanok
  \uses{def:ricci, thm:levi-civita-li}
  $\g \mapsto -2\Ric(\g)$ is $C^n$ at every nondegenerate $\g$, for every $n$.
\end{lemma}

\begin{proof}
  \leanok
  $\Ric$ is a polynomial in $\g$ and $\g^{-1}$. Building every operator inside
  the algebra of continuous linear maps makes smoothness a mechanical chain of
  composition lemmas, with \texttt{contDiffAt\_map\_inverse} at the single
  non-polynomial node. Nondegeneracy of an inner product comes from positive
  definiteness plus a dimension count.
\end{proof}

\begin{theorem}[Short-time existence and uniqueness]
  \label{thm:ricci-lie-shortTime}
  \lean{RicciFlowBlueprint.LeftInvariant.ricciFlow_leftInvariant}
  \leanok
  \uses{lem:ricci-lie-ode}
  Through every left-invariant metric there is a solution of
  $\partial_t \g = -2\Ric(\g)$ near each time $t_0$, and it is unique as a germ
  at $t_0$.
\end{theorem}

\begin{proof}
  \leanok
  Picard--Lindel\"of, via
  \texttt{ContDiffAt.exists\_forall\_mem\_closedBall\_exists\_eq\_forall\_mem\_Ioo\_hasDerivAt}
  for existence and \texttt{ODE\_solution\_unique\_of\_eventually} for
  uniqueness; the 2025 refactor means \texttt{IsPicardLindel\"of} never has to
  be constructed by hand.

  \textbf{Stated as a germ, deliberately.} An earlier version of this blueprint
  claimed a unique \emph{maximal} solution. Mathlib has no maximal-solution
  theory: gluing local solutions along a maximal interval is routine
  mathematics but absent infrastructure. The honest statement is what is
  proved.
\end{proof}

\chapter{Milnor's formulas}
\label{chap:milnor}

Also closed, and also PDE-free. Milnor's 1976 paper computes the curvature of
left-invariant metrics in terms of structure constants; nothing from it had been
formalized. Everything in this chapter is proved.

\textbf{Conventions.} Ours are $\Riem(x,y) = \nab_x\nab_y - \nab_y\nab_x -
\nab_{[x,y]}$ and $\Ric(x,y) = \operatorname{tr}(w \mapsto \Riem(w,x)y)$, as in
Chapter~\ref{chap:curvature}. Milnor flips \emph{both} the curvature sign and
the contraction slot, so his numerical values agree with ours --- the two flips
cancel. Had only one flipped, every eigenvalue below would be silently negated.

\begin{lemma}[The Koszul formula]
  \label{lem:milnor-koszul}
  \lean{RicciFlowBlueprint.LeftInvariant.apply_koszul_apply_milnor}
  \leanok
  \uses{def:koszul-li}
  $\langle \nab_x y, z\rangle = \tfrac12\big(\langle [x,y],z\rangle
     - \langle [y,z],x\rangle + \langle [z,x],y\rangle\big)$.
\end{lemma}

\begin{lemma}[Ricci in structure constants]
  \label{lem:milnor-ricci}
  \lean{RicciFlowBlueprint.LeftInvariant.ricciBilin_structureConstants}
  \leanok
  \uses{lem:milnor-koszul}
  With $\alpha_{ijk} = \langle [e_i,e_j], e_k\rangle$ for an orthonormal basis,
  $\Ric$ expands as a double sum in the $\alpha_{ijk}$ --- in any dimension, for
  any orthonormal basis. This is the computational backbone of the paper.
\end{lemma}

\begin{definition}[Milnor frame]
  \label{def:milnor-frame}
  \lean{RicciFlowBlueprint.LeftInvariant.IsMilnorFrame}
  \leanok
  In dimension three, a basis with
  $[e_2,e_3] = \lambda_1 e_1$, $[e_3,e_1] = \lambda_2 e_2$,
  $[e_1,e_2] = \lambda_3 e_3$.

  \textbf{Existence is proved} --- see Theorem~\ref{thm:milnor-classification}.
\end{definition}

\begin{theorem}[Ricci is diagonal in a Milnor frame]
  \label{thm:milnor-ricci-diag}
  \lean{RicciFlowBlueprint.LeftInvariant.ricciBilin_milnorFrame}
  \leanok
  \uses{lem:milnor-ricci, def:milnor-frame}
  In a Milnor frame $\Ric$ is diagonal, with
  \[ r_i = 2\mu_j\mu_k, \qquad
     \mu_i = \tfrac12(\lambda_j + \lambda_k - \lambda_i), \]
  equivalently $r_i = \tfrac12\big(\lambda_i^2 - (\lambda_j-\lambda_k)^2\big)$.
\end{theorem}

\begin{theorem}[Milnor's classification lemma]
  \label{thm:milnor-classification}
  \lean{RicciFlowBlueprint.LeftInvariant.exists_milnorFrame}
  \leanok
  \uses{def:milnor-frame}
  Every $3$-dimensional unimodular metric Lie algebra admits a Milnor frame.
\end{theorem}

\begin{proof}
  \leanok
  In dimension three the alternating bracket factors as $[x,y] = L(x \times y)$
  through the isomorphism $\Lambda^2\mathfrak{g} \cong \mathfrak{g}$;
  unimodularity says exactly that $L$ is self-adjoint; and the spectral theorem
  supplies an orthonormal eigenbasis, whose eigenvalues are the $\lambda_i$.
  Dimension three is essential --- it is where $\Lambda^2\mathfrak{g} \cong
  \mathfrak{g}$ comes from.
\end{proof}

\begin{corollary}[Ricci is diagonalizable]
  \label{cor:milnor-ricci-diagonalizable}
  \lean{RicciFlowBlueprint.LeftInvariant.exists_milnorFrame_ricci_diagonal}
  \leanok
  \uses{thm:milnor-classification, thm:milnor-ricci-diag}
  Every $3$-dimensional unimodular metric Lie algebra admits a basis in which
  $\Ric$ is diagonal with $r_i = 2\mu_j\mu_k$ --- with no frame assumed.
\end{corollary}

\begin{corollary}[Heisenberg has mixed Ricci signs]
  \label{cor:milnor-heisenberg}
  \lean{RicciFlowBlueprint.LeftInvariant.ricci_heisenberg_mixed_sign}
  \leanok
  \uses{thm:milnor-ricci-diag}
  The Heisenberg algebra ($\lambda_1 = \lambda_2 = 0 \neq \lambda_0$) has
  principal Ricci curvatures
  $(\lambda_0^2/2,\, -\lambda_0^2/2,\, -\lambda_0^2/2)$: strictly mixed signs.
  An instance of Milnor's Theorem 2.4.
\end{corollary}

\begin{theorem}[Milnor's scalar curvature]
  \label{thm:milnor-scalar}
  \lean{RicciFlowBlueprint.LeftInvariant.exists_milnorFrame_scalar}
  \leanok
  \uses{thm:milnor-ricci-diag, thm:milnor-classification, def:scalar}
  Every $3$-dimensional unimodular metric Lie algebra admits an orthonormal
  Milnor frame in which
  \[ \scal = 2(\mu_0\mu_1 + \mu_1\mu_2 + \mu_2\mu_0), \]
  with no frame assumed. At Lie-algebra level raising the index is a single
  composition with $\g^{-1}$, because the metric is carried as a map into the
  dual.

  Sanity check: the round $SU(2)$ has $\lambda_i = 2$, hence $\mu_i = 1$ and
  $\scal = 6$ --- the unit $3$-sphere.
\end{theorem}

\begin{theorem}[The three-dimensional flow, explicitly]
  \label{thm:isenberg-jackson}
  \lean{RicciFlowBlueprint.LeftInvariant.ricciField_milnorFrame_diag}
  \leanok
  \uses{thm:milnor-ricci-diag, thm:ricci-lie-shortTime}
  At a diagonal metric $\operatorname{diag}(A,B,C)$ in a Milnor frame the field
  $\g \mapsto -2\Ric(\g)$ is again diagonal --- the ansatz is invariant --- and
  \[ A' = -\frac{A^2\lambda_0^2}{BC} + \frac{B\lambda_1^2}{C}
          + \frac{C\lambda_2^2}{B} - 2\lambda_1\lambda_2, \]
  cyclically. This is the classical Isenberg--Jackson system: Ricci flow on the
  three-dimensional unimodular Lie groups, written out.

  Proving the $3$-dimensional case for a \emph{diagonal} rather than orthonormal
  metric is what removes the square roots and makes this derivable.

  \textbf{Stated at the level of the field's value} at diagonal metrics. That
  solutions \emph{remain} diagonal needs a second Picard--Lindel\"of pass
  through the positive octant, and is not proved.
\end{theorem}

\chapter{Ricci flow on a closed manifold}
\label{chap:flow}

Everything from here on is blocked on analysis Mathlib does not have. It is
stated so the dependency graph is honest, not because it is close. The point of
writing it out is to see where the road forks and how much of it a given
target actually requires.

\begin{definition}[Ricci flow]
  \label{def:ricci-flow}
  \lean{RicciFlowBlueprint.IsRicciFlowOn}
  \leanok
  \uses{def:ricci}
  A family $\g(t)$ of metrics with $\partial_t \g = -2\Ric(\g)$.

  \textbf{Defined.} Stated per time slice: at each $t$ there is a $C^1$
  torsion-free metric-compatible connection whose Ricci curvature gives the
  rate of change of $\g$. By Theorem \ref{thm:ricci-well-defined} this is a
  statement about $\g(t)$ alone, and Lemma \ref{lem:flow-any-lc} makes that
  explicit. $\forall t\,\exists \nab$ collapses to a family by choice, since no
  regularity in $t$ is demanded of the connection. With the function
  $\g \mapsto \Ric(\g)$ of Definition \ref{def:ricci-of-metric}, the
  existential disappears altogether: Lemma \ref{lem:flow-textbook} is the
  textbook equation.

  The definition is on an interval rather than all of $\R$: flows live on
  intervals, and short-time existence cannot be stated otherwise.
\end{definition}

\begin{lemma}[Stationary flows are Ricci-flat]
  \label{lem:flow-stationary}
  \lean{RicciFlowBlueprint.isRicciFlowAt_const_iff}
  \leanok
  \uses{def:ricci-flow}
  A constant family is a Ricci flow at $t$ if and only if its Levi-Civita
  connection is Ricci-flat.
\end{lemma}

\begin{proof}
  \leanok
  Uniqueness of derivatives against $\partial_t(\text{const}) = 0$, in both
  directions. Proved --- and it is the check that
  Definition~\ref{def:ricci-flow} has content rather than being vacuously
  satisfiable.
\end{proof}

\begin{lemma}[The flow equation is independent of the connection]
  \label{lem:flow-any-lc}
  \lean{RicciFlowBlueprint.isRicciFlowAt_iff_of_isLeviCivita}
  \leanok
  \uses{def:ricci-flow, thm:ricci-well-defined}
  For any $C^1$ Levi-Civita connection $\nab$ of $\g(t)$, the family $\g$ is a
  Ricci flow at $t$ if and only if
  $\partial_t \g(t)(X, Y) = -2\Ric_{\nab}(X, Y)$ for all $C^2$ fields $X, Y$.
\end{lemma}

\begin{proof}
  \leanok
  One direction is the witness. For the other, the witness $\nab'$ of the
  existential and the given $\nab$ have the same $\Ric$ on $C^2$ fields by
  Theorem \ref{thm:ricci-well-defined}.
\end{proof}

\begin{lemma}[Ricci flow as an equation in the metric alone]
  \label{lem:flow-textbook}
  \lean{RicciFlowBlueprint.isRicciFlowOn_iff_ricciOfMetric}
  \leanok
  \uses{def:ricci-flow, def:ricci-of-metric, lem:flow-any-lc}
  $\g$ is a Ricci flow on $s$ if and only if
  $\partial_t \g(t)(X, Y) = -2\Ric(\g(t))(X, Y)$ for all $t \in s$ and all
  $C^2$ fields $X, Y$, with $\Ric(\g(t))$ the Ricci curvature of the metric
  $\g(t)$. No connection is quantified over: this is the textbook equation.
\end{lemma}

\begin{proof}
  \leanok
  Lemma \ref{lem:flow-any-lc} applied to the canonical connection, which is
  $C^1$ by Theorem \ref{thm:levi-civita-smooth}.
\end{proof}

\begin{theorem}[Short-time existence]
  \label{thm:short-time}
  \uses{def:ricci-flow}
  On a closed manifold, Ricci flow has a unique solution on a maximal interval
  $[0,T)$ for any smooth initial metric.

  \textbf{Stated} as \texttt{ricciFlow\_shortTime\_existence}, via
  \texttt{proof\_wanted}: every metric on a closed manifold is the initial
  value of a flow on some $(0,T)$, with coefficients continuous up to $t = 0$.

  \textbf{Blocked.} Requires quasilinear parabolic systems on sections of
  vector bundles (via DeTurck's trick, or Nash--Moser in Hamilton's original),
  and hence Sobolev spaces of bundle sections and elliptic regularity. None of
  this exists in Mathlib. This is the single largest missing prerequisite on
  the whole road, and it is shared by every node below.
\end{theorem}

\begin{definition}[Second covariant derivative]
  \label{def:hessian}
  \lean{CovariantDerivative.hessian}
  \leanok
  \uses{def:curvature}
  $\nab^2_{X,Y} Z = \nab_X \nab_Y Z - \nab_{\nab_X Y} Z$, and for a function
  $\nab^2 f(X,Y) = X(Yf) - (\nab_X Y)f$. Both are tensorial in $X$ and $Y$
  (\texttt{tensorialAt\_hessian\_fst/snd}), so the Hessian of a $C^2$ field at a
  point is a bilinear map on $T_xM$ (\texttt{hessianAt}).
\end{definition}

\begin{lemma}[Ricci identity; symmetry of the Hessian]
  \label{lem:ricci-identity}
  \lean{CovariantDerivative.hessian_sub_hessian_swap}
  \leanok
  \uses{def:hessian, lem:derivation}
  For a torsion-free connection,
  $\nab^2_{X,Y} Z - \nab^2_{Y,X} Z = R(X,Y)Z$, and the Hessian of a $C^2$
  function is symmetric (\texttt{hessianFun\_symm}).
\end{lemma}

\begin{proof}
  \leanok
  Expanding the definitions, the difference is
  $R(X,Y)Z + \nab_{[X,Y]}Z - \nab_{\nab_X Y - \nab_Y X} Z$, and torsion-freeness
  kills the last two terms. For functions the antisymmetric part is
  $[X,Y]f - (\nab_X Y - \nab_Y X)f$, and $[X,Y]f = X(Yf) - Y(Xf)$ is
  Lemma \ref{lem:derivation}.
\end{proof}

\begin{definition}[Laplacian]
  \label{def:laplacian}
  \lean{CovariantDerivative.laplacian, CovariantDerivative.laplacian_eq_sum_frame}
  \leanok
  \uses{def:hessian}
  The rough Laplacian of a $C^2$ vector field is the metric trace of its
  Hessian, $\Delta Z = \sum_i \nab^2_{e_i,e_i} Z$ over an orthonormal basis of
  $T_xM$; it does not depend on the basis (\texttt{laplacian\_eq\_sum}, via
  \texttt{OrthonormalBasis.sum\_apply\_self\_eq}, the statement that
  $\sum_i B(e_i,e_i)$ is basis-independent for any bilinear $B$ --- and note
  that this holds for bilinear maps into any normed space, so no pairing
  against a covector is needed and the trace comes out vector-valued
  directly). \texttt{laplacian\_eq\_sum\_frame} reads the same sum off a
  frame of honest vector fields, which is the form frame computations
  (\texttt{Divergence.lean}) use.

  This is the operator the evolution equations below are written in, and it
  is what supplies the differential inequality at a spatial minimum that the
  abstract maximum principle \ref{thm:max-scalar} takes as its hypothesis;
  the link is the second-derivative test, Lemma \ref{lem:second-derivative-test}.
\end{definition}

\begin{theorem}[The Bochner identity]
  \label{thm:bochner}
  \lean{CovariantDerivative.hessianFun_inner_eq,
    CovariantDerivative.laplacianFun_inner_eq,
    CovariantDerivative.laplacianFun_inner_self_eq,
    CovariantDerivative.two_mul_inner_laplacian_le,
    CovariantDerivative.sum_inner_cov_congr}
  \leanok
  \uses{def:laplacian, def:hessian, thm:levi-civita}
  For a metric connection,
  \[
    \Delta \langle \sigma, \tau\rangle
      = \langle \Delta\sigma, \tau\rangle + \langle \sigma, \Delta\tau\rangle
        + 2\sum_i \langle \nab_{e_i}\sigma, \nab_{e_i}\tau\rangle,
  \]
  and in particular $\Delta |\sigma|^2 = 2\langle \Delta\sigma,\sigma\rangle +
  2\sum_i |\nab_{e_i}\sigma|^2 \ \ge\ 2\langle \Delta\sigma,\sigma\rangle$.

  \textbf{Proved} (\texttt{Bochner.lean}). The pointwise statement behind it is
  the Hessian form,
  \[
    \nab^2\langle\sigma,\tau\rangle(X,Y)
      = \langle \nab^2_{X,Y}\sigma, \tau\rangle
      + \langle \sigma, \nab^2_{X,Y}\tau\rangle
      + \langle \nab_Y\sigma, \nab_X\tau\rangle
      + \langle \nab_X\sigma, \nab_Y\tau\rangle,
  \]
  which is metric compatibility applied twice --- once to differentiate
  $\langle\sigma,\tau\rangle$ along $Y$, once to differentiate the result along
  $X$. The Hessian's correction term $-(\nab_X Y)\langle\sigma,\tau\rangle$ is
  what turns the two iterated derivatives into $\nab^2\sigma$ and
  $\nab^2\tau$; expanded, again by metric compatibility, it supplies exactly
  the two corrections needed.

  The trace is taken over the frame $\Delta$ and $\Delta_{\text{fun}}$ are
  themselves defined on, the canonical extensions of the standard orthonormal
  basis of $T_x\M$, so no frame-independence is needed as an input, and the
  factor $2$ comes from the two cross terms coinciding at $X = Y$.

  \textbf{The gradient term is frame-independent on its own}
  (\texttt{sum\_inner\_cov\_congr}), and this needs neither metric
  compatibility nor the identity above: $(v,w) \mapsto \langle \nab_v\sigma,
  \nab_w\tau\rangle$ is a continuous bilinear form, because $\nab\sigma$ and
  $\nab\tau$ are continuous linear maps at each point, so its trace is
  basis-independent for the usual reason.

  The inequality $\Delta|\sigma|^2 \ge 2\langle\Delta\sigma,\sigma\rangle$ is
  the form the maximum principle \ref{thm:max-scalar} consumes: the discarded
  term is a sum of squares.
\end{theorem}

\begin{lemma}[$\nab \Rm$ is tensorial in all four slots]
  \label{lem:cov-curvature-direction}
  \lean{CovariantDerivative.covCurvature_smul_dir,
    CovariantDerivative.covCurvature_add_dir,
    CovariantDerivative.cov_smul_dir,
    CovariantDerivative.contMDiff_curvature,
    CovariantDerivative.covCurvature_antisymm,
    CovariantDerivative.covCurvature_smul_snd,
    CovariantDerivative.covCurvature_add_snd,
    CovariantDerivative.covCurvature_smul_thd,
    CovariantDerivative.covCurvature_add_thd,
    CovariantDerivative.covCurvature_smul_fth,
    CovariantDerivative.covCurvature_add_fth}
  \leanok
  \uses{lem:bianchi-second, def:curvature, lem:curvature-pointwise}
  $(\nab_X \Riem)(Y,Z)W$ is $C^\infty(M)$-linear in each of $X$, $Y$, $Z$ and
  $W$ separately: replacing any one of them by $fV$ multiplies the value at $x$
  by $f(x)$, and replacing it by a sum splits the value.

  \textbf{Proved} (\texttt{CurvatureDeriv.lean}). Three different mechanisms,
  not one repeated.

  \emph{The direction slot} is the cheap one. Each of the four terms of
  $\nab\Riem$ picks up exactly one factor of $f(x)$, for the same reason:
  $\nab_{fX}$ is $f\nab_X$ on the nose, and the last three terms feed that into
  a slot of $\Riem$ in which $\Riem$ is already tensorial --- the first two
  cheaply, the third by Lemma \ref{lem:curvature-pointwise}.

  \emph{The $Y$ and $Z$ slots} are a cancellation instead. Under $Y \mapsto fY$
  the \emph{first} term picks up a Leibniz term,
  $\nab_X(f\,\Riem(Y,Z)W) = f\nab_X(\Riem(Y,Z)W) + (Xf)\Riem(Y,Z)W$, cancelled
  by the matching term from $\Riem(\nab_X(fY),Z)W$ --- so the direction slot was
  the easy one precisely because it has none. Applying Leibniz needs $y \mapsto
  \Riem(Y,Z)W(y)$ to be a \emph{differentiable section}, which
  \texttt{covCurvature} does not require to be well defined. That is
  \texttt{contMDiff\_curvature}, also proved here: a $C^3$ triple has a $C^1$
  curvature section, the two $\nab\nab$ terms by
  \texttt{contMDiff\_cov\_apply} twice and the bracket term by the same fed
  Mathlib's \texttt{ContMDiffAt.mlieBracket\_vectorField}. Once $Y$ is done,
  \texttt{covCurvature\_antisymm} hands over $Z$ for nothing.

  \emph{The $W$ slot} is the same cancellation with the other pairing:
  $\Riem(Y,Z)(fW) = f\,\Riem(Y,Z)W$ as \emph{sections}, so the leading $\nab_X$
  contributes $(Xf)\Riem(Y,Z)W$, and the last term contributes the matching one
  through $\nab_X(fW) = f\nab_X W + (Xf)W$. It costs one derivative more of $f$
  than the others, and that is not slack: the second Leibniz term is
  $\Riem(Y,Z)\big((Xf)W\big)$, and tensoriality in $\Riem$'s third slot wants
  its coefficient $Xf$ to be $C^2$, hence $f \in C^3$. Stating this needed
  \texttt{curvature\_smul\_third'}, Lemma \ref{lem:curvature-pointwise}'s
  third-slot linearity restated asking of the first two slots only the $C^1$
  regularity its proof actually uses --- the fields landing in them here are
  $\nab_X Y$ and $\nab_X Z$, which are one derivative worse than $Y$ and $Z$.

  With all four slots, the \emph{divergence}
  $\operatorname{div}\Rm(Y,Z,W) = \sum_i \langle (\nab_{e_i}\Riem)(Y,Z)W, e_i
  \rangle$ is well defined --- the trace is over the direction slot, and
  tensoriality is what makes the sum independent of the orthonormal frame. The
  divergence is what the contracted second Bianchi identity, and hence $\Delta
  \scal$, are about.
\end{lemma}

\begin{lemma}[$\nab \Rm$ is pointwise in all four slots]
  \label{lem:cov-curvature-pointwise}
  \lean{CovariantDerivative.covCurvature_congr_snd,
    CovariantDerivative.covCurvature_congr_thd,
    CovariantDerivative.covCurvature_congr_fth,
    CovariantDerivative.covCurvature_zero_snd,
    CovariantDerivative.covCurvature_sum_smul_snd,
    CovariantDerivative.covCurvature_congr_snd_of_eventuallyEq,
    CovariantDerivative.covCurvature_zero_fth,
    CovariantDerivative.covCurvature_sum_smul_fth,
    CovariantDerivative.covCurvature_congr_fth_of_eventuallyEq,
    CovariantDerivative.contMDiff_curvature_two,
    CovariantDerivative.contMDiff_covCurvature}
  \leanok
  \uses{lem:cov-curvature-direction, lem:curvature-pointwise, lem:global-extension}
  $(\nab_X\Riem)(Y,Z)W$ at $x$ depends only on the four values $X(x)$, $Y(x)$,
  $Z(x)$, $W(x)$. Together with Lemma \ref{lem:cov-curvature-direction} this
  makes $\nab\Riem$ a genuine tensor, not merely an operator on fields.

  \textbf{Proved} (\texttt{CurvatureDeriv.lean}). The direction slot was free
  (Definition \ref{def:div-curvature}). The other three are the
  frame-and-globalise argument of Lemma \ref{lem:curvature-pointwise} one level
  up: expand the field in a local frame, globalise the frame fields and the
  coefficients, apply tensoriality, and read off that the coefficients at $x$
  depend only on the value at $x$. Locality --- the input that argument needs
  --- is cheap in each case, because every term of $\nab\Riem$ other than the
  leading one is $\Riem$ in a slot where $\Riem$ is already tensorial, hence
  pointwise; only $\nab_X(\Riem(Y,Z)W)$ needs the germ.

  Two asymmetries are worth recording. The $Z$ slot is free from the $Y$ slot by
  \texttt{covCurvature\_antisymm}, as with tensoriality. And the $W$ slot runs
  at one derivative higher: \texttt{covCurvature\_smul\_fth} asks its
  coefficient for three derivatives, so the frame fields and the coefficient
  functions must be globalised at $C^3$ rather than $C^2$, and the frame comes
  from \texttt{IsOrthonormalFrameOn}\ $\ldots 3$.

  Also proved here, as the first step towards $\Delta\Rm$:
  \texttt{contMDiff\_curvature\_two} (the curvature of $C^4$ fields is a $C^2$
  section) and \texttt{contMDiff\_covCurvature} ($\nab\Riem$ of $C^3$ fields,
  $C^4$ in the last slot, is a $C^1$ section) --- what is needed to
  differentiate $\nab\Riem$ a second time.
\end{lemma}

\begin{definition}[$\nab^2 \Rm$ and the Laplacian of the curvature tensor]
  \label{def:cov2-curvature}
  \lean{CovariantDerivative.cov2Curvature,
    CovariantDerivative.cov2Curvature_congr_dir,
    CovariantDerivative.cov2Curvature_smul_dir,
    CovariantDerivative.cov2Curvature_add_dir,
    CovariantDerivative.cov2Curvature_smul_snd,
    CovariantDerivative.cov2Curvature_add_snd,
    CovariantDerivative.cov2Curvature_congr_snd,
    CovariantDerivative.cov2CurvatureAt,
    CovariantDerivative.cov2CurvatureBilin,
    CovariantDerivative.curvatureLaplacian,
    CovariantDerivative.curvatureLaplacian_eq_sum_basis,
    CovariantDerivative.curvatureLaplacian_eq_sum_frame}
  \leanok
  \uses{lem:cov-curvature-pointwise, lem:cov-curvature-direction, lem:global-extension}
  \[
    (\nab^2_{X,Y}\Riem)(A,B)C
      = \nab_X\big((\nab_Y\Riem)(A,B)C\big)
        - (\nab_{\nab_X Y}\Riem)(A,B)C
        - (\nab_Y\Riem)(\nab_X A,B)C
        - (\nab_Y\Riem)(A,\nab_X B)C
        - (\nab_Y\Riem)(A,B)(\nab_X C),
  \]
  one correction for each of $\nab\Riem$'s four slots, and
  \[
    \Delta\Rm(A,B)C \;=\; \sum_i (\nab^2_{e_i,e_i}\Riem)(A,B)C .
  \]

  \textbf{Defined and proved frame-independent} (\texttt{CurvatureLaplacian.lean}).

  The outer derivative slot is free, and this is where the four-slot work of
  Lemma \ref{lem:cov-curvature-pointwise} is cashed in: every term is either a
  continuous linear map evaluated at $X(x)$, or $\nab\Riem$ in a slot where it
  is pointwise, applied to a field whose value at $x$ is determined by $X(x)$.
  So \texttt{cov2Curvature\_congr\_dir} needs no frame argument at all.

  The inner slot is the familiar cancellation --- the Leibniz term of
  $\nab_X\big(f\,(\nab_Y\Riem)(A,B)C\big)$ is matched by the one
  $\nab_X(fY) = f\nab_X Y + (Xf)Y$ produces in the second term --- and its
  pointwise statement is the frame-and-globalise argument once more, at $C^3$,
  since the coefficient must be $C^3$ for $(Xf)$ to be $C^2$.

  With both derivative slots pointwise and bilinear, $\nab^2\Riem$ is a
  continuous bilinear map $T_x\M \times T_x\M \to T_x\M$ and $\Delta\Rm$ is its
  trace; frame-independence is then \texttt{OrthonormalBasis.sum\_apply\_self\_eq},
  which traces bilinear maps into any normed space, so the result is
  vector-valued directly with no pairing against a covector.

  Regularity: $X, Y, A, B$ are $C^3$ and $C$ is $C^4$. The last slot is the
  expensive one at every level of this tower --- $C$ needs four derivatives
  because \texttt{covCurvature\_congr\_fth} needs $C^3$ coefficients because
  \texttt{curvature\_smul\_third} needs $C^2$ ones.

  \textbf{What this is for.} $\Delta\Rm$ is the right-hand side of Hamilton's
  evolution equation $\partial_t \Rm = \Delta\Rm + Q(\Rm)$, which is what
  carries Hamilton--Ivey pinching from the ODE (Theorem
  \ref{thm:hamilton-ivey}) to the flow. The geometry side of that equation is
  now in place; the analytic side, $\partial_t\Rm$, still lives on the model
  space (\texttt{CurvatureVariation.lean}).
\end{definition}

\begin{definition}[The divergence of the curvature tensor]
  \label{def:div-curvature}
  \lean{CovariantDerivative.covCurvature_congr_dir,
    CovariantDerivative.covCurvatureAt,
    CovariantDerivative.covCurvatureEndo,
    CovariantDerivative.divCurvature,
    CovariantDerivative.divCurvature_eq_sum_basis,
    CovariantDerivative.divCurvature_eq_sum_frame}
  \leanok
  \uses{lem:cov-curvature-direction, lem:global-extension}
  \[
    \operatorname{div}\Rm(Y,Z,W)(x)
      \;=\; \sum_i \big\langle (\nab_{e_i}\Riem)(Y,Z)W,\ e_i \big\rangle,
  \]
  the trace of $\nab\Riem$ over its direction slot against its fourth.

  \textbf{Defined and proved frame-independent} (\texttt{Divergence.lean}). Two
  facts make this well posed, and neither is free.

  First, $(\nab_X\Riem)(Y,Z)W$ at $x$ depends only on $X(x)$
  (\texttt{covCurvature\_congr\_dir}). This costs no new frame argument, unlike
  Lemma \ref{lem:curvature-pointwise}: each of the four terms is separately
  pointwise in $X$. The first is a continuous linear map evaluated at $X(x)$;
  the next two are $\Riem$ in a slot where it is tensorial, applied to
  $\nab_X Y$ and $\nab_X Z$, whose values at $x$ are $\nab_{X(x)}Y$ and
  $\nab_{X(x)}Z$; the last is $\Riem$'s third slot applied to $\nab_X W$, which
  is Lemma \ref{lem:curvature-pointwise} itself.

  Second, it is $C^\infty(\M)$-linear there
  (Lemma \ref{lem:cov-curvature-direction}). So
  $v \mapsto (\nab_v\Riem)(Y,Z)W$ is a genuine endomorphism of $T_x\M$ ---
  \texttt{covCurvatureEndo}, built on globally $C^2$ extensions of $v$ as
  \texttt{ricciAt} is --- and $\operatorname{div}\Rm$ is its trace.

  \textbf{Frame-independence is then algebraic}, not a computation: it is
  \texttt{LinearMap.trace\_eq\_sum\_inner}, valid for any orthonormal basis of
  $T_x\M$ (\texttt{divCurvature\_eq\_sum\_basis}), and
  \texttt{divCurvature\_eq\_sum\_frame} reads it off any $C^2$ frame whose
  values at $x$ are orthonormal. Note the contrast with the scalar curvature:
  there the object traced is a bilinear form and the trace is metric; here it
  is an endomorphism and no metric enters the trace at all --- the metric
  appears only in writing it as $\sum_i \langle \cdot, e_i\rangle$.

  This is the object the contracted second Bianchi identity is about, and
  through that identity the object $\Delta \scal$ is about.
\end{definition}

\begin{lemma}[The traced second Bianchi identity]
  \label{lem:traced-bianchi}
  \lean{CovariantDerivative.divCurvature_eq_sub_sum}
  \leanok
  \uses{def:div-curvature, lem:bianchi-second}
  \[
    \operatorname{div}\Rm(Y,Z,W)
      \;=\; \sum_i \big\langle (\nab_Y\Riem)(e_i,Z)W, e_i\big\rangle
          \;-\; \sum_i \big\langle (\nab_Z\Riem)(e_i,Y)W, e_i\big\rangle .
  \]

  \textbf{Proved} (\texttt{Divergence.lean}). Take the second Bianchi identity
  with $e_i$ in the direction slot,
  \[
    (\nab_{e_i}\Riem)(Y,Z)W + (\nab_Y\Riem)(Z,e_i)W + (\nab_Z\Riem)(e_i,Y)W = 0,
  \]
  pair with $e_i$, and sum. The middle term is turned round by antisymmetry of
  $\nab\Riem$ in its second and third slots
  (\texttt{covCurvature\_antisymm}). Nothing is differentiated along the frame:
  the identity is pointwise in $e_i$ and is only summed, so no derivative of
  $e_i$ appears and no orthonormality is used beyond what
  Definition \ref{def:div-curvature} already needs.

  \textbf{What is left to reach $\operatorname{div}\Ric = \tfrac12 d\scal$} is a
  single further identity: the two sums on the right are $(\nab_Y\Ric)(Z,W)$ and
  $(\nab_Z\Ric)(Y,W)$. That is the trace commuting with $\nab$ for the curvature
  endomorphism --- the analogue for $\Riem$ of Lemma
  \ref{lem:metric-trace-cov}, and like it a statement about a frame that is not
  parallel, so the coefficients $\langle \nab_X e_i, e_j\rangle$ enter and are
  killed by antisymmetry. It is not proved yet.
\end{lemma}

\begin{theorem}[The contracted second Bianchi identity, first contraction]
  \label{thm:contracted-bianchi}
  \lean{CovariantDerivative.covRicci,
    CovariantDerivative.curvatureBilinFst,
    CovariantDerivative.sum_inner_covCurvature_eq,
    CovariantDerivative.divCurvature_eq_covRicci_sub,
    CovariantDerivative.ricci_eq_sum_inner_frame}
  \leanok
  \uses{lem:traced-bianchi, lem:metric-trace-cov, def:div-curvature}
  For a Levi-Civita connection,
  \[
    \operatorname{div}\Rm(Y,Z,W) \;=\; (\nab_Y\Ric)(Z,W) - (\nab_Z\Ric)(Y,W).
  \]

  \textbf{Proved} (\texttt{Divergence.lean}). Lemma \ref{lem:traced-bianchi}
  already gives this with the right-hand side written as two frame sums
  $\sum_i \langle (\nab_X\Riem)(e_i,Z)W, e_i\rangle$. What is proved here is
  that each such sum \emph{is} $(\nab_X\Ric)(Z,W)$ ---
  \texttt{sum\_inner\_covCurvature\_eq}, the statement that the first-slot trace
  commutes with $\nab$.

  \textbf{Why that is not immediate.} $\Ric(Z,W)$ is a frame sum only where the
  frame is orthonormal, so differentiating it in $X$ means differentiating the
  frame too, and a $C^2$ orthonormal frame is not parallel. Expanding
  $X\big(\sum_i \langle \Riem(e_i,Z)W, e_i\rangle\big)$ by metric compatibility
  produces, besides the wanted terms, $\sum_i \langle \Riem(e_i,Z)W, \nab_X
  e_i\rangle$; and $(\nab_X\Riem)(e_i,Z)W$ carries the matching $-\Riem(\nab_X
  e_i, Z)W$. Their sum is
  \[
    \sum_i \Big[\,\beta(\nab_X e_i,\, e_i) + \beta(e_i,\, \nab_X e_i)\,\Big],
    \qquad \beta(a,b) := \langle \Riem(a,Z)W,\, b\rangle,
  \]
  which vanishes because the coefficients $\langle \nab_X e_i, e_j\rangle$ are
  antisymmetric (metric compatibility plus local constancy of $\langle e_i,
  e_j\rangle$) while the bracket is symmetric in $i,j$. That is exactly the
  mechanism of Lemma \ref{lem:metric-trace-cov}, and the same Lean lemma
  \texttt{sum\_bilin\_of\_antisymm} does the work; what is new is that $\beta$
  has to be produced as a genuine continuous bilinear form, which is
  \texttt{curvatureBilinFst}, available only because $\Riem$ is tensorial in its
  first slot.

  Both hypotheses on the connection are used, and for different things:
  torsion-freeness for the second Bianchi identity, metric compatibility for the
  trace to commute with $\nab$.

  \textbf{What is left} is the second contraction, $\operatorname{div}\Ric =
  \tfrac12\, d\scal$: trace this identity again over $Y$ and $W$. That needs the
  pair symmetry of $\Riem$ carried through $\nab$, and Lemma
  \ref{lem:metric-trace-cov} to turn $\sum_j (\nab_Z\Ric)(e_j,e_j)$ into
  $Z(\scal)$.
\end{theorem}

\begin{theorem}[The contracted second Bianchi identity, second contraction]
  \label{thm:div-ricci}
  \lean{CovariantDerivative.inner_covCurvature_expand,
    CovariantDerivative.inner_covCurvature_pair_symm,
    CovariantDerivative.covBilin_ricciForm_eq_covRicci,
    CovariantDerivative.covRicci_symm,
    CovariantDerivative.two_mul_sum_covRicci_eq_mvfderiv_scalarCurvatureAt}
  \leanok
  \uses{thm:contracted-bianchi, lem:metric-trace-cov, lem:curvature-symmetries,
    thm:ricci-symm}
  For a Levi-Civita connection,
  $\operatorname{div}\Ric = \tfrac12\, d\scal$; in the frame form the Lean
  statement takes,
  \[
    2 \sum_i (\nab_{e_i}\Ric)(e_i, Y) \;=\; Y(\scal).
  \]

  \textbf{Proved} (\texttt{Divergence.lean}). Contract Theorem
  \ref{thm:contracted-bianchi} over its remaining two slots: sum
  \[
    \operatorname{div}\Rm(Y,e_j,e_j) = (\nab_Y\Ric)(e_j,e_j) -
    (\nab_{e_j}\Ric)(Y,e_j)
  \]
  over $j$. On the right the first sum is $Y(\scal)$ by Lemma
  \ref{lem:metric-trace-cov}, and the second is $\operatorname{div}\Ric(Y)$
  because $\nab\Ric$ is symmetric --- which it is because $\Ric$ is (Lemma
  \ref{thm:ricci-symm}), so this step too needs the connection to be
  metric. On the left, pair symmetry of $\nab\Riem$ turns $\langle
  (\nab_{e_i}\Riem)(Y,e_j)e_j, e_i\rangle$ into $\langle
  (\nab_{e_i}\Riem)(e_j,e_i)Y, e_j\rangle$, whose sum over $j$ is
  $(\nab_{e_i}\Ric)(e_i,Y)$ by Theorem \ref{thm:contracted-bianchi}'s trace
  lemma; summing over $i$ gives $\operatorname{div}\Ric(Y)$ back. So
  $\operatorname{div}\Ric(Y) = Y(\scal) - \operatorname{div}\Ric(Y)$.

  \textbf{Pair symmetry of $\nab\Riem$ is not inherited for free.} What Lemma
  \ref{lem:curvature-symmetries} gives is a symmetry of the $(0,4)$ tensor
  $\langle\Riem(A,B)C,D\rangle$, and $\nab\Riem$ is defined as an operator on
  fields, not as the derivative of that tensor. The bridge is
  \texttt{inner\_covCurvature\_expand}: for a \emph{metric} connection the
  leading $\nab_X$ can be moved onto the pairing, so
  \[
    \langle (\nab_X\Riem)(A,B)C, D\rangle
      = X\langle \Riem(A,B)C,D\rangle - \sum_{\text{4 slots}} \langle
      \Riem(\ldots,\nab_X\cdot,\ldots),\cdot\rangle - \langle \Riem(A,B)C,
      \nab_X D\rangle,
  \]
  five corrections rather than four --- the fifth is the metric one. Every term
  on the right is the $(0,4)$ tensor evaluated somewhere, so the pointwise
  symmetry transfers, and the two sides' corrections are the same five numbers
  in a different order.

  With this, roadmap item 2 is closed. The next link in the chain is
  $\operatorname{tr}_g \partial_t \Ric = \Delta \scal$ under the flow; combined
  with \texttt{RicciVariation.lean}'s $\partial_t \scal = \operatorname{tr}_g
  \partial_t\Ric + 2|\Ric|^2$ that gives $\partial_t \scal = \Delta \scal +
  2|\Ric|^2$.
\end{theorem}

\begin{lemma}[$\nab h$ is a tensor]
  \label{lem:cov-bilin-tensor}
  \lean{CovariantDerivative.covBilin_smul_dir,
    CovariantDerivative.covBilin_add_dir,
    CovariantDerivative.covBilin_congr_dir,
    CovariantDerivative.covBilin_smul_snd,
    CovariantDerivative.covBilin_add_snd,
    CovariantDerivative.covBilin_smul_thd,
    CovariantDerivative.covBilin_add_thd,
    CovariantDerivative.covBilinDir,
    CovariantDerivative.covBilinForm,
    CovariantDerivative.covBilinForm_apply,
    CovariantDerivative.sum_covBilin_congr}
  \leanok
  \uses{def:laplacian, lem:metric-trace-cov}
  Let $h$ be a bilinear form field on $TM$ and $\nab$ a connection. Then
  $(\nab_X h)(Y,Z) = X(h(Y,Z)) - h(\nab_X Y, Z) - h(Y, \nab_X Z)$ is
  \emph{tensorial and pointwise in all three slots}: a continuous trilinear
  form on $T_xM$, not merely an operation on fields. In particular
  $\nab_X h$ is an honest bilinear form $\texttt{covBilinForm}$, and the
  direction is an honest linear form $\texttt{covBilinDir}$.

  \textbf{Proved} (\texttt{BilinDeriv.lean}). All three slots are cheap,
  unlike the four slots of $\nab\Rm$ (Lemma \ref{lem:cov-curvature-direction}).
  The \emph{direction} slot needs \emph{nothing}: each of the three terms is a
  continuous linear map --- $\mathrm{d}(h(Y,Z))_x$, $(\nab Y)_x$, $(\nab Z)_x$
  --- applied to $X(x)$, so linearity and pointwiseness hold by construction,
  with no differentiability hypothesis whatever. The $Y$ and $Z$ slots are the
  usual Leibniz cancellation: replacing $Y$ by $fY$ produces $(Xf)\,h(Y,Z)$
  from the derivative term, and $-h(\nab_X(fY),Z)$ produces $-(Xf)\,h(Y,Z)$,
  which cancels it; only $\mathrm{MDifferentiableAt}$ of the fields is used,
  as in the Hessian's second slot, so no frame-and-globalise argument of the
  kind $\Rm$'s third slot needs (Lemma \ref{lem:curvature-pointwise}) appears.
  Hence \texttt{TensorialAt.mkHom} and \texttt{mkHom}$_2$ apply directly.

  The one hypothesis carried is \texttt{IsMDiffBilinAt}: that $y \mapsto h_y(U_y,
  V_y)$ is differentiable at $x$ for differentiable $U$, $V$. It is what makes
  the Leibniz rule available and is implied by $h$ being $C^1$ into the space
  of bilinear forms.

  \textbf{Consequence.} $\operatorname{tr}_\g(\nab_X h)$ is frame-independent
  on its own (\texttt{sum\_covBilin\_congr}), by the same trace-of-a-bilinear-map
  argument that Definition \ref{def:ricci-form} uses for $\scal$; so both sides
  of Lemma \ref{lem:metric-trace-cov} are frame-independent, not just equal for
  each frame. This is the object the connection Laplacian $\Delta_\g h$ and the
  double divergence $\operatorname{div}\operatorname{div} h$ are traces of, and
  those are what $\operatorname{tr}_\g(\partial_t \Ric) = \Delta \scal$ needs.
\end{lemma}

\begin{proof}
  \leanok
  Direction: by construction, as above. Second slot: with
  $\nab_X(fY) = f\nab_X Y + (Xf) Y$ and $h(fY,Z) = f\,h(Y,Z)$,
  \[
    (\nab_X h)(fY,Z) = (Xf)h(Y,Z) + f\,X(h(Y,Z)) - f\,h(\nab_XY,Z)
      - (Xf)h(Y,Z) - f\,h(Y,\nab_XZ) = f\,(\nab_X h)(Y,Z),
  \]
  and additivity is the same computation without the coefficient. The third
  slot is identical with the roles of $Y$ and $Z$ exchanged.
\end{proof}

\begin{definition}[$\nab^2 h$ and the connection Laplacian $\Delta_\g h$]
  \label{def:bilin-laplacian}
  \lean{CovariantDerivative.cov2Bilin,
    CovariantDerivative.cov2Bilin_congr_dir,
    CovariantDerivative.cov2Bilin_smul_dir,
    CovariantDerivative.cov2Bilin_add_dir,
    CovariantDerivative.cov2Bilin_smul_snd,
    CovariantDerivative.cov2Bilin_add_snd,
    CovariantDerivative.cov2Bilin_congr_snd,
    CovariantDerivative.cov2BilinAt,
    CovariantDerivative.cov2BilinForm,
    CovariantDerivative.laplacianBilin,
    CovariantDerivative.laplacianBilin_eq_sum_basis,
    CovariantDerivative.laplacianBilin_eq_sum_frame}
  \leanok
  \uses{lem:cov-bilin-tensor, def:laplacian, def:cov2-curvature}
  For a bilinear form field $h$,
  \[
    (\nab^2_{W,X} h)(Y,Z) = W\big((\nab_X h)(Y,Z)\big) - (\nab_{\nab_W X} h)(Y,Z)
      - (\nab_X h)(\nab_W Y, Z) - (\nab_X h)(Y, \nab_W Z)
  \]
  --- one correction per slot of $\nab h$, so \emph{three}, against the five of
  $\nab^2\Rm$ (Definition \ref{def:cov2-curvature}). The \textbf{connection
  Laplacian} $\Delta_\g h$ is its metric trace in the two derivative slots,
  $(\Delta_\g h)(Y,Z) = \sum_i (\nab^2_{e_i,e_i} h)(Y,Z)$.

  \textbf{Proved} (\texttt{BilinLaplacian.lean}). $\nab^2 h$ is tensorial and
  pointwise in both derivative slots, so $\Delta_\g h$ is well defined and
  frame-independent. The two slots behave differently, as they do for
  $\nab^2\Rm$:

  \begin{itemize}
    \item The \emph{outer} direction $W$ needs \textbf{no frame argument}. It is
      Lemma \ref{lem:cov-bilin-tensor} cashed in: the leading term is a
      continuous linear map at $W(x)$, and each of the three corrections is
      $\nab h$ in a slot it is already pointwise in, applied to a field whose
      value at $x$ is determined by $W(x)$.
    \item The \emph{inner} direction $X$ is the Leibniz cancellation --- and the
      cancellation is cheap, because every step other than the leading one uses
      only the \emph{direction} slot of $\nab h$, which carries no hypotheses at
      all. Pointwiseness there is the frame-and-globalise argument, but it runs
      at $C^1$: \texttt{cov2Bilin\_smul\_snd} asks its coefficient for one
      derivative, against the three that $\nab^2\Rm$ asks.
  \end{itemize}

  The one hypothesis beyond Lemma \ref{lem:cov-bilin-tensor}'s is
  \texttt{IsMDiffCovBilinAt}: that $y \mapsto (\nab_V h)(Y,Z)$ is differentiable
  at $x$ for $C^1$ fields $V$. Merely-differentiable $V$ cannot supply this ---
  the leading term of $\nab h$ differentiates $h(Y,Z)$ in the direction $V$, so
  it sees the germ of $V$ --- which is exactly why the inner slot needs the
  frame argument and the outer one does not.

  This is what $\operatorname{tr}_\g(\partial_t\Ric) = \Delta \scal$ is
  assembled from. Tracing the first variation of $\Rm$ (Lemma
  \ref{lem:variation-curvature}) twice against the Koszul formula for
  $\partial_t\nab$ gives, with \emph{no} curvature terms,
  \[
    \operatorname{tr}_\g(\partial_t \Ric)
      = \operatorname{div}\operatorname{div} h - \Delta(\operatorname{tr}_\g h),
  \]
  and under the flow $h = -2\Ric$, so $\operatorname{tr}_\g h = -2\scal$ and
  $\operatorname{div}\operatorname{div} h = -2\operatorname{div}\operatorname{div}\Ric
  = -\Delta \scal$ by the second contraction (Theorem \ref{thm:div-ricci}),
  giving $2\Delta\scal - \Delta\scal = \Delta\scal$.
\end{definition}

\begin{definition}[$\nab\omega$, $\operatorname{div}\omega$, and the double divergence of a bilinear form]
  \label{def:one-form-divergence}
  \lean{CovariantDerivative.covOneForm,
    CovariantDerivative.covOneForm_smul_dir,
    CovariantDerivative.covOneForm_smul_snd,
    CovariantDerivative.covOneFormAt,
    CovariantDerivative.divOneForm,
    CovariantDerivative.divOneForm_eq_sum_basis,
    CovariantDerivative.divOneForm_eq_sum_frame,
    CovariantDerivative.covBilinFstSnd,
    CovariantDerivative.divBilin,
    CovariantDerivative.divBilin_eq_sum_basis,
    CovariantDerivative.divBilin_eq_sum_frame,
    CovariantDerivative.divBilinForm,
    CovariantDerivative.divDivBilin,
    CovariantDerivative.covOneForm_divBilinOneForm_eq_sum,
    CovariantDerivative.divDivBilin_eq_sum}
  \leanok
  \uses{lem:cov-bilin-tensor, def:bilin-laplacian, def:laplacian}
  For a one-form field $\omega$, $(\nab_X \omega)(Y) = X(\omega(Y)) -
  \omega(\nab_X Y)$, and $\operatorname{div}\omega = \sum_i (\nab_{e_i}
  \omega)(e_i)$. For a bilinear form field $h$, $(\operatorname{div} h)(Z) =
  \sum_i (\nab_{e_i} h)(e_i, Z)$, and the \textbf{double divergence} is
  $\operatorname{div}\operatorname{div} h = \operatorname{div}
  (\operatorname{div} h)$.

  \textbf{Proved} (\texttt{OneForm.lean}). The point is that
  $\operatorname{div}\operatorname{div} h$ --- the trace of $\nab^2 h$ pairing
  slots $1\&3$ and $2\&4$ --- needs \emph{no four-linear packaging of $\nab^2h$
  at all}. It factors into two cheap traces:

  \begin{itemize}
    \item $\operatorname{div} h$ is the trace of the $(0,3)$-tensor $\nab h$
      (Lemma \ref{lem:cov-bilin-tensor}) in its first two slots, which is one
      $\texttt{mkHom}_2$ over (direction, second slot): the direction slot is
      tensorial with no hypotheses, and the second is the Leibniz cancellation.
      It is tensorial in its remaining slot as a finite sum of third-slot
      tensorialities, so it is an honest one-form field.
    \item $\nab\omega$ has \emph{two} terms, against $\nab h$'s three and
      $\nab\Rm$'s four, so \emph{both} of its slots are free of any frame
      argument: the direction is a continuous linear map by construction, and
      the $Y$ slot is the plain Leibniz cancellation on merely differentiable
      fields. So $\texttt{TensorialAt.mkHom}_2$ applies directly and
      $\operatorname{div}\omega$ is frame-independent for free.
  \end{itemize}

  \textbf{The two definitions agree.} $\operatorname{div}\operatorname{div} h$
  is defined here as an \emph{iterated} divergence, while the evolution of the
  scalar curvature produces the \emph{double trace} of $\nab^2h$. They are equal:
  \[
    (\nab_W(\operatorname{div} h))(Z) = \sum_i (\nab^2_{W,e_i}h)(e_i, Z),
    \qquad
    \operatorname{div}\operatorname{div} h
      = \sum_{j}\sum_i (\nab^2_{e_j,e_i}h)(e_i, e_j)
  \]
  (\texttt{covOneForm\_divBilinOneForm\_eq\_sum},
  \texttt{divDivBilin\_eq\_sum}) --- \emph{the divergence commutes with
  $\nab$}. This is Lemma \ref{lem:metric-trace-cov} applied to the bilinear
  form $(v,w) \mapsto (\nab_v h)(w,Z)$: differentiating the frame sum produces
  $\nab^2h$ \emph{plus} exactly the term $(\operatorname{div} h)(\nab_W Z)$
  that $\nab_W$ of a one-form subtracts, so the two cancel and no curvature
  term survives.

  For $\omega = \mathrm{d}f$ these are the Hessian and Laplacian of a function
  (Definition \ref{def:laplacian}); the general case is what the flow needs.
  With $\Delta_\g h$ from Definition \ref{def:bilin-laplacian}, both traces in
  \[
    \operatorname{tr}_\g(\partial_t \Ric)
      = \operatorname{div}\operatorname{div} h - \Delta(\operatorname{tr}_\g h)
  \]
  now exist, and under the flow $h = -2\Ric$ the second contraction of the
  second Bianchi identity (Theorem \ref{thm:div-ricci}) turns the right-hand
  side into $2\Delta\scal - \Delta\scal = \Delta\scal$.
\end{definition}

\begin{lemma}[The metric trace commutes with $\nab$]
  \label{lem:metric-trace-cov}
  \lean{CovariantDerivative.mvfderiv_sum_eq_sum_covBilin,
    CovariantDerivative.mvfderiv_sum_eq_sum_covBilin_of_frame,
    CovariantDerivative.inner_cov_antisymm,
    CovariantDerivative.sum_bilin_of_antisymm,
    CovariantDerivative.mvfderiv_scalarCurvatureAt_eq_sum_covBilin}
  \leanok
  \uses{def:laplacian, thm:levi-civita, def:ricci-form}
  Let $\nab$ be a metric connection, $B$ a bilinear form field on $TM$ and
  $e_1, \dots, e_n$ a local orthonormal frame near $x$. Then
  \[
    X\Big(\sum_i B(e_i, e_i)\Big) \;=\; \sum_i (\nab_X B)(e_i, e_i),
  \]
  where $(\nab_X B)(Y,Z) = X(B(Y,Z)) - B(\nab_X Y, Z) - B(Y, \nab_X Z)$. Since
  $\sum_i B(e_i,e_i)$ is the metric trace $\operatorname{tr}_\g B$ for any
  orthonormal frame, this is $X(\operatorname{tr}_\g B) =
  \operatorname{tr}_\g(\nab_X B)$.

  \textbf{Proved} (\texttt{TraceCov.lean}). This is the gate for every
  Laplacian identity below: the contracted second Bianchi identity, $\partial_t
  R = \Delta R + 2|\Ric|^2$ inside Lemma \ref{lem:variation-scalar}, and
  $\partial_t \Rm = \Delta \Rm + Q$ (Lemma \ref{lem:evolution-rm}).

  $B$ need not be symmetric, and the frame need not be parallel --- which is
  the point, since a parallel frame exists only along a curve. Two auxiliary
  results are stated separately and are of independent use:
  \texttt{inner\_cov\_antisymm}, that the connection coefficients of an
  orthonormal frame are antisymmetric, and
  \texttt{sum\_bilin\_of\_antisymm}, the purely algebraic cancellation.
  A local orthonormal frame is supplied by \texttt{OrthonormalFrame.lean}, and
  \texttt{exists\_orthonormalBasis\_of\_isOrthonormalFrameOn} turns it into
  an orthonormal basis of each fibre over its base set.

  \textbf{First consequence}, and the reason for the node:
  $X(\scal) = \operatorname{tr}_\g(\nab_X \Ric)$
  (\texttt{mvfderiv\_scalarCurvatureAt\_eq\_sum\_covBilin}), taking $B$ to be
  the Ricci form of Definition \ref{def:ricci-form}. The scalar curvature is the
  frame sum at \emph{every} point of the frame's base set, so the two sides are
  germ-equal at $x$ and the trace lemma applies verbatim. Its one remaining
  hypothesis is genuine regularity --- $y \mapsto \Ric(e_i,e_i)(y)$ must be
  differentiable, which asks the connection for one derivative more than
  $\texttt{ContMDiffCovariantDerivative}\ 1$ gives.
\end{lemma}

\begin{proof}
  \leanok
  Expanding the definition of $\nab_X B$, the claim is that the $2n^2$
  correction terms cancel:
  \[
    \sum_i \big[ B(\nab_X e_i, e_i) + B(e_i, \nab_X e_i) \big] = 0 .
  \]
  Write $\nab_X e_i = \sum_j a_{ij} e_j$ with $a_{ij} = \langle \nab_X e_i,
  e_j \rangle$. Because $\langle e_i, e_j \rangle$ is constant near $x$, its
  derivative along $X$ vanishes, and metric compatibility
  $X\langle e_i, e_j\rangle = \langle \nab_X e_i, e_j\rangle + \langle e_i,
  \nab_X e_j\rangle$ forces $a_{ij} = -a_{ji}$. The sum above is then
  $\sum_{i,j} a_{ij}\big[B(e_j,e_i) + B(e_i,e_j)\big]$, in which the bracket is
  symmetric under $i \leftrightarrow j$ while $a$ is antisymmetric; swapping
  the two summation indices shows the sum is its own negative. Differentiating
  the trace term by term needs $\mathrm{d}$ of a finite sum
  (\texttt{mvfderiv\_fun\_sum}), and the vanishing of the derivative of a
  locally constant function needs that $\mathrm{d}$ depends only on the germ
  (\texttt{Filter.EventuallyEq.mvfderiv\_eq}); both are missing from Mathlib
  and are proved here.
\end{proof}

\begin{lemma}[Second-derivative test]
  \label{lem:second-derivative-test}
  \lean{RicciFlowBlueprint.laplacianFun_nonneg_of_isLocalMin}
  \leanok
  \uses{def:laplacian}
  At a local minimum $x_0$ of a $C^2$ function $f$ on a boundaryless
  manifold, $\nab f(x_0) = 0$, $\nab^2 f(X,X)(x_0) \ge 0$ for every vector
  field $X$ differentiable at $x_0$, and $\Delta f(x_0) \ge 0$; for any
  covariant derivative, not only the Levi-Civita connection.

  \textbf{Proved}, first on the model space
  (\texttt{hessianFun\_nonneg\_of\_isLocalMin\_model}), then on a manifold by
  transport through \texttt{extChartAt}
  (\texttt{hessianFun\_nonneg\_of\_isLocalMin}), since at a critical point
  the Hessian is connection-independent. With it, the abstract maximum
  principle \ref{thm:max-scalar} becomes the classical one for
  $\partial_t u = \Delta u + F(u)$.
\end{lemma}

\begin{proof}
  \leanok
  On the line: if $g'' (t_0) < 0$ at a local minimum with $g'(t_0) = 0$,
  then $g' < 0$ just to the right of $t_0$ and the mean value theorem gives
  $g(t_1) < g(t_0)$ there (\texttt{deriv2\_nonneg\_of\_isLocalMin}). In a
  normed space, restrict to the line $x_0 + tv$ to get
  $D^2 f(x_0)(v,v) \ge 0$ (\texttt{fderiv2\_nonneg\_of\_isLocalMin}). For the
  Hessian, $\nab^2 f(X,X) = X(Xf) - (\nab_X X) f$; the product rule gives
  $X(Xf) = D^2f(X,X) + Df(DX\cdot X)$, and $Df(x_0) = 0$ kills both the second
  term and the connection term
  (\texttt{fderiv\_fderiv\_apply\_nonneg\_of\_isLocalMin}). On a manifold,
  write $g = f \circ (\mathrm{extChartAt}\, x_0)^{-1}$: $g$ has a local
  minimum at the chart image of $x_0$ since the inverse chart is continuous,
  so $Dg = 0$ there; $X(Xf)$ at $x_0$ is the chart-side expression
  $D(Dg \cdot \tilde X)(\tilde X)$ with $\tilde X$ the pullback of $X$ (the
  same transport as Lemma \ref{lem:derivation}), and the vector-space case
  applies. Boundarylessness is what makes $g$ genuinely $C^2$ at an interior
  point of the model rather than $C^2$ within its range. The Laplacian is a
  sum of such terms over an orthonormal basis, each extended to a
  differentiable section.
\end{proof}

\begin{lemma}[First variation of the Levi-Civita connection]
  \label{lem:variation-connection}
  \lean{RicciFlowBlueprint.exists_hasDerivAt_leviCivitaOfMetricE_inner_eq}
  \leanok
  \uses{thm:levi-civita-unique, def:ricci-flow}
  Along a family of metrics $\g(t)$ with $\partial_t \g = h$, the Levi-Civita
  connections $\nab^t$ satisfy
  \[
    \g\bigl(\partial_t \nab_X Y,\, Z\bigr)
      = \tfrac12\bigl[(\nab_X h)(Y,Z) + (\nab_Y h)(Z,X) - (\nab_Z h)(X,Y)\bigr],
  \]
  with $\nab = \nab^{t_0}$ on the right and
  $(\nab_X h)(Y,Z) = X(h(Y,Z)) - h(\nab_X Y, Z) - h(Y, \nab_X Z)$
  (\texttt{CovariantDerivative.covBilin}).

  \textbf{Proved}, including the differentiability of $t \mapsto \nab^t_X Y(x)$,
  with one hypothesis stated in exactly the form used: the space derivatives
  $X(\g(Y,Z))$ commute with $\partial_t$ for every test field $Z$
  differentiable at $x$ (three permuted forms), a joint-regularity statement
  about $(t,x) \mapsto \g(t)(x)$ which the pointwise-in-$t$ Definition
  \ref{def:ricci-flow} does not impose. Differentiability of
  $t \mapsto \nab^t_X Y(x)$ comes from the musical isomorphism:
  $\nab^t_X Y = (\g_t)^{\sharp\,-1}(K_t)$ with $K_t = \g_t(\nab^t_X Y, \cdot)$
  the Koszul functional; $(\g_t)^\sharp$ is invertible by positive
  definiteness (\texttt{isInvertible\_innerE}), $K_t$ is differentiable
  coordinatewise by the differentiated Koszul formula, and
  \texttt{ContinuousLinearMap.inverse} is smooth at invertible points
  (\texttt{exists\_hasDerivAt\_leviCivitaOfMetricE}). The scalar form
  (\texttt{hasDerivAt\_inner\_leviCivitaOfMetric}) needs only the commutation
  hypothesis for the one field $Z$ at hand.
\end{lemma}

\begin{proof}
  \leanok
  By the Koszul formula, $\g(\nab_X Y, Z)$ is an explicit expression in the
  metric alone, so $\partial_t$ of it is the same expression in $h$
  (\texttt{hasDerivAt\_inner\_leviCivitaOfMetric}). The Koszul combination of a
  symmetric bilinear form $h$, rewritten through a torsion-free connection, is
  $(\nab_X h)(Y,Z) + (\nab_Y h)(Z,X) - (\nab_Z h)(X,Y) + 2h(\nab_X Y, Z)$
  (\texttt{koszul\_bilin\_eq}: for $h = \g$ and $\nab$ compatible the first three
  terms vanish and this is Koszul itself). The product rule
  $\partial_t[\g(\nab_X Y, Z)] = h(\nab_X Y, Z) + \g(\partial_t \nab_X Y, Z)$
  absorbs the last term.

  Time derivatives of the metric are typed on the model space $E = T_xM$:
  \texttt{TangentSpace} carries no norm of its own (the norm comes from the
  metric), so \texttt{HasDerivAt} into $T_xM \to T_xM \to \R$ cannot be stated,
  while into $E \to E \to \R$ it can, and in finite dimensions the notion does
  not depend on the norm.
\end{proof}

\begin{lemma}[First variation of the curvature]
  \label{lem:variation-curvature}
  \lean{RicciFlowBlueprint.hasDerivAt_curvatureE_leviCivitaOfMetric}
  \leanok
  \uses{def:curvature, lem:variation-connection}
  Along a family of metrics $\g(t)$ with $\partial_t \g = h$, the curvature
  tensors $\Rm^t$ of the Levi-Civita connections satisfy
  \[
    \partial_t \Rm(X,Y)Z = (\nab_X \dot A)(Y,Z) - (\nab_Y \dot A)(X,Z),
  \]
  with $\nab = \nab^{t_0}$ and $\dot A = \partial_t \nab$ the tensor of
  Lemma \ref{lem:variation-connection},
  $\g(\dot A(v,w), Z) = \tfrac12[(\nab_v h)(w,Z) + (\nab_w h)(Z,v) - (\nab_Z h)(v,w)]$
  (\texttt{inner\_derivDifferenceE\_eq}). With $h = -2\Ric$ this is
  $\partial_t \Rm$ under the flow, before the second Bianchi identity turns
  $\nab^2 \Ric$ into $\Delta \Rm + Q(\Rm)$ (Lemma \ref{lem:evolution-rm}).

  \textbf{Proved}, in two layers. The algebraic layer
  (\texttt{curvature\_eq\_add\_covEnd}) is for any two connections
  $\nab' = \nab + A$ with $\nab$ torsion-free:
  \[
    R'(X,Y)Z = R(X,Y)Z + (\nab_X A)(Y,Z) - (\nab_Y A)(X,Z)
      + A(X, A(Y,Z)) - A(Y, A(X,Z)).
  \]
  The analytic layer (\texttt{hasDerivAt\_curvatureE}) differentiates this
  along a family $\nab^t = \nab + A^t$ with $A^{t_0} = 0$: the quadratic
  terms have zero derivative at $t_0$. Two hypotheses are stated in the form
  used: the commutation of Lemma \ref{lem:variation-connection}, now for all
  fields differentiable at the point (\texttt{CommutesWithMvfderiv}), and the
  commutation of $\partial_t$ with $\nab_X$ on the sections $A^t(Y,Z)$ and
  $A^t(X,Z)$ at $x$. Both are joint regularity in $(t,x)$.
\end{lemma}

\begin{proof}
  \leanok
  The difference of two connections is a tensor, Mathlib's
  \texttt{CovariantDerivative.difference}; expanding
  $\nab'_X \nab'_Y Z - \nab'_Y \nab'_X Z - \nab'_{[X,Y]} Z$ with
  $\nab' = \nab + A$ and collecting, the terms $A(\nab_X Y, Z) - A(\nab_Y X, Z)$
  combine with $-A([X,Y], Z)$ by torsion-freeness, and what remains is the
  displayed identity. Along the family, $t \mapsto A^t(y)$ is differentiable as a
  bilinear map because it is differentiable on every pair of vectors
  (\texttt{exists\_hasDerivAt\_clm\_of\_apply}, finite dimensions) and each value
  $A^t(v,w)(y) = \nab^t_v \tilde w - \nab_v \tilde w$ is differentiable by
  Lemma \ref{lem:variation-connection}, $\tilde w$ the extension of $w$ constant
  in a trivialisation. The derivative $\dot A$ is \texttt{deriv}, so the
  statement carries no existential.
\end{proof}

\begin{lemma}[Variation of the metric trace]
  \label{lem:metric-trace-variation}
  \lean{RicciFlowBlueprint.hasDerivAt_metricTraceE}
  \leanok
  \uses{def:scalar, lem:variation-connection}
  For a metric $\g$ on $T_xM$ and a bilinear form $B$, write
  $\operatorname{tr}_\g B = \sum_i B(e_i,e_i)$ over a $\g$-orthonormal basis,
  equivalently the trace of the endomorphism $\g^{\sharp\,-1} B^\flat$. Along a
  family of metrics with $\partial_t \g = h$ and a family of forms with
  $\partial_t B = \dot B$,
  \[
    \partial_t \operatorname{tr}_{\g_t} B_t
      = \operatorname{tr}_\g \dot B - \langle h, B\rangle_\g,
    \qquad
    \langle h, B\rangle_\g = \sum_{i,j} B(e_i,e_j)\, h(e_j,e_i).
  \]
  This is what turns Lemma \ref{lem:variation-curvature} into the variations of
  $\Ric = \operatorname{tr}_\g \Rm$ and $R = \operatorname{tr}_\g \Ric$; under
  the flow $h = -2\Ric$, and the correction term is the $2|\Ric|^2$ of
  $\partial_t R = \Delta R + 2|\Ric|^2$.

  \textbf{Proved}, at a point, on the model space $E = T_xM$
  (\texttt{MetricTrace.lean}). The trace is defined as
  $\operatorname{tr}(\g^{\sharp\,-1} \circ B^\flat)$
  (\texttt{metricTraceE}) and shown equal to the sum over any orthonormal basis
  (\texttt{metricTraceE\_eq\_sum}); the pairing is the trace of
  $\g^{\sharp\,-1} h^\flat \g^{\sharp\,-1} B^\flat$
  (\texttt{metricTraceE\_comp\_sharpE\_eq\_sum}).
\end{lemma}

\begin{proof}
  \leanok
  The inverse of the metric is differentiable, since inversion of continuous
  linear maps is smooth at invertible points, and its derivative is found by
  differentiating $\g^\sharp \circ \g^{\sharp\,-1} = \mathrm{id}$:
  $\partial_t \g^{\sharp\,-1} = -\g^{\sharp\,-1} h^\flat \g^{\sharp\,-1}$
  (\texttt{hasDerivAt\_inverse\_innerE}). The product rule for
  $\g^{\sharp\,-1} \circ B^\flat$ and linearity of the trace give the formula.
\end{proof}

\begin{lemma}[First variations of the Ricci and scalar curvatures]
  \label{lem:variation-scalar}
  \lean{RicciFlowBlueprint.hasDerivAt_ricciOfMetric,
    RicciFlowBlueprint.hasDerivAt_scalarCurvatureOfMetric'_of_isRicciFlowAt,
    RicciFlowBlueprint.ricciFormOfMetric,
    RicciFlowBlueprint.ricciFormOfMetric_apply_field,
    RicciFlowBlueprint.innerE_deriv_eq_of_isRicciFlowAt',
    RicciFlowBlueprint.scalarCurvatureOfMetricAt,
    RicciFlowBlueprint.scalarCurvatureOfMetricAt_eq_metricTraceE,
    RicciFlowBlueprint.hasDerivAt_ricciFormOfMetric,
    RicciFlowBlueprint.hasDerivAt_scalarCurvatureOfMetricAt,
    RicciFlowBlueprint.hasDerivAt_scalarCurvatureOfMetricAt_of_isRicciFlowAt}
  \leanok
  \uses{lem:variation-curvature, lem:metric-trace-variation, def:ricci-of-metric, def:scalar, def:ricci-flow}
  Along a family of metrics $\g(t)$ with $\partial_t \g = h$,
  \[
    \partial_t \Ric(X,Y) = \operatorname{tr}\bigl(\partial_t [v \mapsto \Rm(v,X)Y]\bigr),
    \qquad
    \partial_t R = \operatorname{tr}_\g(\partial_t \Ric) - \langle h, \Ric\rangle_\g,
  \]
  and under the Ricci flow $h = -2\Ric$, so
  $\partial_t R = \operatorname{tr}_\g(\partial_t \Ric) + 2|\Ric|^2$.

  \textbf{Proved} (\texttt{RicciVariation.lean}). Ricci is the trace of the
  endomorphism $v \mapsto \Rm(v,X)Y$ of $T_xM$, so its variation is the trace
  of the variation of Lemma \ref{lem:variation-curvature} and no term from the
  metric appears; this is on any manifold. Scalar curvature is the
  \emph{metric} trace of the Ricci form, and Lemma
  \ref{lem:metric-trace-variation} supplies the correction $-\langle h,
  \Ric\rangle_\g$, which is $2|\Ric|^2_\g = 2\sum_{i,j} \Ric(e_i,e_j)\Ric(e_j,e_i)$
  under the flow.

  \textbf{The scalar statement now holds on any manifold.} It used to be
  confined to the model space $M = E$, because it was phrased through
  \texttt{ricciE}, the Ricci form built on \emph{constant} fields, and those
  are global sections only when $M = E$. The restriction is removed by working
  with the tensors of Definition \ref{def:ricci-form} instead:
  \texttt{ricciFormOfMetric} is the Ricci form of $\g$ on any manifold ---
  Definition \ref{def:ricci-form} for the Levi-Civita connection of $\g$,
  built on the global $C^2$ extensions of Lemma \ref{lem:global-extension} ---
  and \texttt{scalarCurvatureOfMetricAt} is $\operatorname{scal}$ likewise.
  Three things then transfer verbatim.
  \texttt{innerE\_deriv\_eq\_of\_isRicciFlowAt'} carries the flow equation
  $\partial_t \g = -2\Ric$ across, by instantiating the flow on those
  extensions and using uniqueness of derivatives;
  \texttt{hasDerivAt\_ricciFormOfMetric} differentiates the Ricci form
  coordinatewise, each value $\Ric_t(v,w)(x)$ being $\Ric_t(V,W)(x)$ for any
  $C^2$ fields with $V(x) = v$, $W(x) = w$; and
  \texttt{scalarCurvatureOfMetricAt\_eq\_metricTraceE} identifies
  $\operatorname{scal}$ with $\operatorname{tr}_\g \Ric$, so Lemma
  \ref{lem:metric-trace-variation} applies unchanged. The conclusions are
  \texttt{hasDerivAt\_scalarCurvatureOfMetricAt} and, under the flow,
  \texttt{hasDerivAt\_scalarCurvatureOfMetricAt\_of\_isRicciFlowAt}. The
  model-space statements are kept: they are the same theorems where Definition
  \ref{def:scalar} was first phrased, and they need no extension lemma.

  So the only ingredient still missing from $\partial_t R = \Delta R +
  2|\Ric|^2$ is $\operatorname{tr}_\g(\partial_t\Ric) = \Delta R$. That
  rests on the contracted second Bianchi identity, which is now proved in both
  contractions (Theorems \ref{thm:contracted-bianchi} and
  \ref{thm:div-ricci}) rather than pending.
\end{lemma}

\begin{proof}
  \leanok
  $\Ric_t(X,Y)(x) = \operatorname{tr} T_t$ with $T_t v = \Rm_t(\bar v, X)Y(x)$,
  $\bar v$ the extension of $v$ constant in a trivialisation; $t \mapsto T_t$
  is differentiable as a linear map because each $T_t v$ is, by Lemma
  \ref{lem:variation-curvature} on the field $\bar v$, and the trace is linear.
  On the model space the Ricci form $\Ric_t(v,w) = \Ric_t(\bar v, \bar w)(x)$
  is bilinear by tensoriality in both slots (the second slot from the
  curvature's third), differentiable in $t$ coordinatewise, and $R_t =
  \operatorname{tr}_{\g_t} \Ric_t$; Lemma \ref{lem:metric-trace-variation}
  differentiates the trace. Under the flow, $\partial_t \g(v,w) = -2\Ric(v,w)$
  by definition of the flow and uniqueness of derivatives, and
  $\langle -2\Ric, \Ric\rangle_\g = -2|\Ric|^2_\g$.
\end{proof}

\begin{lemma}[Evolution of the curvature tensor]
  \label{lem:evolution-rm}
  \uses{def:ricci-flow, def:curvature, def:laplacian, lem:variation-connection, lem:variation-curvature, lem:bianchi-second, lem:metric-trace-variation, lem:variation-scalar, lem:metric-trace-cov}
  Under the flow, $\partial_t \Rm = \Delta \Rm + Q(\Rm)$, where $Q$ is
  quadratic in the curvature operator.

  Obtained cleanly in an evolving orthonormal frame (Uhlenbeck's trick), which
  removes the terms coming from the time-dependence of the metric on the
  bundle. Formally this is a computation, not an estimate, but it is a long one
  and it is what every subsequent estimate is applied to.

  \textbf{First half done.} Lemma \ref{lem:variation-curvature} gives
  $\partial_t \Rm$ as $\nab \dot A$, second covariant derivatives of
  $h = -2\Ric$. The second Bianchi identity is Lemma \ref{lem:bianchi-second}.
  What remains is its trace against the metric, which rewrites the
  antisymmetrised $\nab^2 \Ric$ as $\Delta \Rm$ plus terms quadratic in
  $\Rm$.
\end{lemma}

\begin{theorem}[Scalar maximum principle]
  \label{thm:max-scalar}
  \lean{RicciFlowBlueprint.MaximumPrinciple.le_of_deriv_ge_at_min, RicciFlowBlueprint.MaximumPrinciple.le_of_laplacian}
  \leanok
  \uses{def:laplacian, lem:second-derivative-test}
  A solution of $\partial_t u = \Delta u + \langle X, \nab u\rangle + F(u)$ on
  a closed manifold, with $F$ Lipschitz, is bounded below by the solution of
  the associated ODE $\dot\phi = F(\phi)$ with $\phi(0) \le \min u(0,\cdot)$;
  and above by it when $\phi(0) \ge \max u(0,\cdot)$.

  \textbf{Proved, abstractly.} The Laplacian enters the classical proof at
  one point only: at a spatial minimum $x_0$ of $u(t,\cdot)$, $\nab u = 0$ and
  $\Delta u \ge 0$, hence $\partial_t u(t,x_0) \ge F(u(t,x_0))$. The
  formalised statement takes exactly that as its hypothesis, on an arbitrary
  compact space: if at every spatial minimum $x_0$ of $u(t,\cdot)$ one has
  $F(u(t,x_0)) \le \partial_t u(t,x_0)$, then $\phi \le u$. No Laplacian and
  no manifold appear; the classical statement is this theorem plus the
  second-derivative test at a minimum, and that is the manifold version below.

  \textbf{And on a manifold.} With the Laplacian of Definition
  \ref{def:laplacian} and the second-derivative test
  (Lemma \ref{lem:second-derivative-test}), the hypothesis at a minimum is
  discharged: \texttt{le\_of\_laplacian} is the statement for
  $\partial_t u = \Delta u + F(u)$ on a closed Riemannian manifold with
  $u(t,\cdot)$ of class $C^2$, for any connection. The gradient term
  $\langle X, \nab u\rangle$ is not yet included: it needs $\nab u = 0$ at a
  minimum, which is the first-derivative test on a manifold.
\end{theorem}

\begin{proof}
  \leanok
  Perturb to $\psi_\epsilon = \phi - \epsilon e^{(2K+1)t}$, a strict
  subsolution by the Lipschitz bound: $\dot\psi_\epsilon < F(\psi_\epsilon)$.
  If $u \ge \psi_\epsilon$ failed, the set of times at which it fails is closed
  (it is the projection of a compact subset of $[0,T] \times M$), so there is a
  first time $t_0 > 0$ and a point $x_0$ with $u(t_0,x_0) = \psi_\epsilon(t_0)$.
  Then $x_0$ is a spatial minimum of $u(t_0,\cdot)$, the left derivative of
  $u(\cdot,x_0) - \psi_\epsilon$ at $t_0$ is $\le 0$, but the differential
  inequality makes it $\ge \epsilon (K+1) e^{(2K+1)t_0} > 0$. Let
  $\epsilon \to 0$. The upper bound is the same argument applied to $-u$.
\end{proof}

\begin{theorem}[Tensor maximum principle]
  \label{thm:max-tensor}
  \lean{RicciFlowBlueprint.MaximumPrinciple.mem_of_deriv_le_at_max, RicciFlowBlueprint.MaximumPrinciple.mem_of_laplacian}
  \leanok
  \uses{thm:max-scalar, def:laplacian, lem:second-derivative-test}
  Hamilton's maximum principle for symmetric tensors: a closed convex subset of
  the bundle preserved by the ODE $\dot{\Rm} = Q(\Rm)$ is preserved by the PDE.

  This is the workhorse. Every curvature-positivity statement in the subject is
  an instance of it, including the one Hamilton needs in dimension three.

  \textbf{Proved, abstractly.} Let $V$ be a complete real inner product space,
  $K \subseteq V$ closed, convex and nonempty, $M$ compact, $u : [0,T] \times M
  \to V$ jointly continuous with time derivative $\partial_t u$, and $F$
  Lipschitz. Two hypotheses are taken in exactly the form the classical proof
  uses them. First, at every spatial maximum $x_0$ of $x \mapsto \langle n,
  u(t,x)\rangle$ one has $\langle n, \partial_t u\rangle \le \langle n,
  F(u)\rangle$ at $(t,x_0)$; on a manifold this is $\partial_t u = \Delta u +
  F(u)$ together with $\Delta \langle n, u\rangle \le 0$ at a maximum. Second,
  $K$ is preserved by the ODE in Nagumo's form: $\langle n, F(p)\rangle \le 0$
  for every $p \in K$ and every outward normal $n$ at $p$, which
  \texttt{subtangential\_of\_invariant} derives from invariance along solution
  curves (the easy direction of Nagumo's theorem). Then $u(0,\cdot) \in K$
  implies $u(t,\cdot) \in K$ for all $t \in [0,T]$. No manifold, no Laplacian
  and no tensor bundle appear; the trivial bundle with fibre $V$ is the case
  Hamilton uses in dimension three, and the general vector-bundle version is the
  same argument in a parallel frame.

  \textbf{And on a manifold.} \texttt{mem\_of\_laplacian} is the statement on
  a closed Riemannian manifold for the trivial bundle $M \times V$, the equation
  $\partial_t u = \Delta u + F(u)$ taken componentwise, $\langle n,
  \partial_t u\rangle = \Delta\langle n, u\rangle + \langle n, F(u)\rangle$
  for every fixed $n \in V$, with $u(t,\cdot)$ of class $C^2$. A maximum of
  $\langle n, u\rangle$ is a minimum of $\langle -n, u\rangle$, so the
  second-derivative test applies unchanged. Hamilton's setting, a bundle with
  an evolving fibre metric, needs the Laplacian on sections together with
  Uhlenbeck's trick.
\end{theorem}

\begin{proof}
  \leanok
  Same skeleton as the proof of \ref{thm:max-scalar}, with a supporting
  half-space of the convex set in place of the half-line $[\phi(t),\infty)$.
  The principle itself does not depend on the evolution equation
  \ref{lem:evolution-rm}; that enters only when it is applied, in Lemma
  \ref{lem:pinching}. Perturb: the claim is $\operatorname{dist}(u(t,x), K) < \epsilon
  e^{(2L+1)t}$. If it fails, the set of times where it fails is closed, so
  there is a first time $t_0 > 0$ and a point $x_0$ where it does; let $p$ be
  the nearest point of $K$ to $u_0 = u(t_0,x_0)$ and $n = u_0 - p$. Since $K$
  lies in the half-space $\langle n, \cdot - p\rangle \le 0$, the function
  $\langle n, u(t,\cdot) - p\rangle$ is bounded by $|n| \cdot
  \operatorname{dist}(u(t,\cdot), K)$, so $x_0$ is a spatial maximum of it at
  time $t_0$, and $t \mapsto \langle n, u(t,x_0) - p\rangle - |n| \epsilon
  e^{(2L+1)t}$ is negative before $t_0$ and zero at $t_0$, whence its left
  derivative is $\ge 0$. The two inequalities give $(2L+1)|n|^2 \le \langle n,
  F(u_0)\rangle \le \langle n, F(p)\rangle + L|n|^2 \le L|n|^2$, absurd. Let
  $\epsilon \to 0$ and use closedness of $K$.
\end{proof}

\begin{theorem}[Shi's local derivative estimates]
  \label{thm:shi}
  \uses{lem:evolution-rm}
  A bound on $|\Rm|$ on a parabolic cylinder gives interior bounds on every
  covariant derivative $|\nab^k \Rm|$, with constants depending only on $k$,
  the dimension, and the geometry of the cylinder.
\end{theorem}

\begin{theorem}[Long-time existence criterion]
  \label{thm:long-time}
  \uses{thm:short-time, thm:shi}
  The flow extends past $T$ if and only if $|\Rm|$ stays bounded on $[0,T)$.
  Singularities are exactly curvature blow-ups.
\end{theorem}

\begin{definition}[Generalized Ricci flow]
  \label{def:generalized-flow}
  \uses{def:ricci-flow}
  The flow recast as a single geometric object: a space-time with a time
  function and a horizontal metric satisfying the flow equation, rather than a
  one-parameter family of metrics on a fixed manifold.
  \emph{(Morgan--Tian 3.8.)}

  Introduced long before surgery and load-bearing for it. Because the
  time-slices are allowed to change topology, surgery becomes an equation on one
  object rather than a procedure applied between objects --- which is exactly
  what makes it a plausible formalization target at all.
\end{definition}

\chapter{Hamilton's theorem}
\label{chap:hamilton}

This is the first genuine milestone, and the last one reachable by
maximum-principle estimates alone. Positive Ricci curvature is a strong enough
hypothesis that the flow never forms a singularity in an interesting way: it
shrinks to a round point.

\begin{lemma}[Hamilton's curvature ODE in dimension three]
  \label{lem:pinching-ode}
  \lean{RicciFlowBlueprint.Pinching.IsCurvatureODE.pinching_antitone}
  \leanok
  In dimension three the curvature operator has three eigenvalues
  $\lambda \ge \mu \ge \nu$ --- twice the sectional curvatures, normalised so
  that the scalar curvature is the trace $\scal = \lambda+\mu+\nu$ and the
  Ricci eigenvalues are
  $\tfrac12(\mu+\nu), \tfrac12(\lambda+\nu), \tfrac12(\lambda+\mu)$ --- and the reaction term
  $\Rm^2 + \Rm^\#$ of the evolution equation \ref{lem:evolution-rm} is
  diagonal in the same frame. The associated ODE is
  \[
    \dot\lambda = \lambda^2 + \mu\nu,\quad
    \dot\mu = \mu^2 + \lambda\nu,\quad
    \dot\nu = \nu^2 + \lambda\mu .
  \]
  Along its solutions on $[0,T]$:
  \begin{enumerate}
  \item the ordering $\lambda \ge \mu \ge \nu$ is preserved
    (\texttt{le\_preserved\_lm}, \texttt{le\_preserved\_mn});
  \item positive Ricci curvature, $\mu + \nu > 0$, is preserved
    (\texttt{ricci\_pos\_preserved});
  \item $\lambda \le C(\mu+\nu)$ is preserved for every $C \ge 1/2$
    (\texttt{bound\_preserved});
  \item given (1)--(3) at time $0$, for $0 \le \delta$ with
    $\delta(2C+1) \le 1$ the ratio $(\lambda-\nu)/(\mu+\nu)^{1-\delta}$ is
    nonincreasing (\texttt{pinching\_antitone}).
  \end{enumerate}
  Item (4) is the pinching estimate of Hamilton's paper (Theorem 10.1) for the
  ODE: the traceless part of the curvature is dominated by a smaller power of
  the scalar curvature, so wherever the curvature blows up it becomes
  constant-sectional to leading order.
\end{lemma}

\begin{proof}
  \leanok
  Each quantity $f$ satisfies $\dot f = a f + (\text{a term of known sign})$
  for a continuous $a$: $(\mu-\lambda)\dot{} = (\mu-\lambda)(\lambda+\mu-\nu)$
  exactly; $(\mu+\nu)\dot{} = \lambda(\mu+\nu) + \mu^2+\nu^2$;
  $(\lambda - C(\mu+\nu))\dot{} = \lambda(\lambda - C(\mu+\nu)) +
  (\mu\nu - C(\mu^2+\nu^2))$ with the bracket $\le 0$ because
  $\mu\nu \le (\mu^2+\nu^2)/2$. So $f e^{-\int a}$ is monotone, which is the
  linear Grönwall comparison \texttt{nonpos\_of\_deriv\_le\_mul}. For (4),
  the logarithmic derivative of the ratio is
  $\delta\lambda + (\nu-\mu) - (1-\delta)(\mu^2+\nu^2)/(\mu+\nu)$, and with
  $\lambda \le C(\mu+\nu)$, $\nu \le \mu$ and $\mu^2+\nu^2 \ge (\mu+\nu)^2/2$
  this is at most $(\mu+\nu)\,(\delta(2C+1)-1)/2 \le 0$.

  No manifold appears: this is the algebraic heart of Hamilton's theorem,
  and everything else in the chapter is the machinery that transfers it from
  the ODE to the flow.
\end{proof}

\begin{lemma}[Pinching in dimension three]
  \label{lem:pinching}
  \uses{thm:max-tensor, lem:pinching-ode, lem:evolution-rm}
  In dimension three positive Ricci curvature is preserved by the flow, and
  the eigenvalues pinch together relative to their size: the traceless part
  of $\Ric$ is dominated by $\scal^{1-\delta}$ for some $\delta > 0$.

  This is Lemma \ref{lem:pinching-ode} transported from the ODE to the PDE by
  the tensor maximum principle \ref{thm:max-tensor}: the sets
  $\{\mu+\nu \ge \epsilon\,\scal\}$, $\{\lambda \le C(\mu+\nu)\}$ and
  $\{\lambda-\nu \le C'(\mu+\nu)^{1-\delta}\}$ are closed, convex, and
  ODE-invariant. Convexity: the largest eigenvalue is convex and the smallest
  concave in the curvature operator, so $\lambda - \nu$ and $\lambda$ are
  convex, $\mu + \nu = \scal - \lambda$ is concave, and $t \mapsto
  t^{1-\delta}$ is concave increasing; each set is a sublevel set of a convex
  function against a concave one.

  \textbf{What is missing} is more than a transfer. In order: the evolution
  equation \ref{lem:evolution-rm} itself; the maximum principle on the bundle
  of curvature operators with the metric evolving on the fibres (Uhlenbeck's
  trick), where \ref{thm:max-tensor} is proved for a fixed fibre; the
  convexity of the eigenvalue-defined sets, which needs the variational
  characterisation of eigenvalues; and Nagumo's condition for these sets from
  the ODE invariance of \ref{lem:pinching-ode}, which is stated for
  eigenvalues and must be lifted to the operator ODE by equivariance.

  Special to dimension three, where the Weyl tensor vanishes and $\Rm$ is
  determined by $\Ric$. The higher-dimensional analogues (Hamilton 1986 in
  dimension four, then B\"ohm--Wilking) are substantially harder and are not on
  this road.
\end{lemma}

\begin{lemma}[Gradient estimate for the scalar curvature]
  \label{lem:gradient-scal}
  \uses{thm:max-tensor, lem:pinching}
  $|\nab \scal| \le \eta\, \scal^{3/2}$ up to lower-order terms, so the scalar
  curvature is comparable at every pair of points as the flow shrinks.
\end{lemma}

\begin{theorem}[Convergence of the normalized flow]
  \label{thm:normalized-convergence}
  \uses{lem:pinching, lem:gradient-scal, thm:long-time, def:sectional}
  Rescaling to fixed volume, the normalized flow exists for all time and
  converges exponentially in every $C^k$ to a metric of constant positive
  sectional curvature.
\end{theorem}

\begin{theorem}[Hamilton 1982]
  \label{thm:hamilton}
  \uses{thm:normalized-convergence, lem:bianchi-first, def:sectional, thm:ricci-well-defined, def:ricci-of-metric, lem:global-extension}
  A closed $3$-manifold admitting a metric of strictly positive Ricci curvature
  admits a metric of constant positive sectional curvature, and is therefore
  diffeomorphic to a spherical space form $S^3/\Gamma$.

  \emph{J. Differential Geometry} \textbf{17} (1982), 255--306.

  \textbf{Stated formally} in \texttt{RicciFlowBlueprint/Hamilton.lean} as
  \texttt{hamilton\_1982}, via \texttt{proof\_wanted} --- so the statement is
  elaborated and type-checked, with no \texttt{sorry} and no added axiom. The
  hypothesis and conclusion are
  \texttt{AdmitsPositiveRicciMetric} and \texttt{AdmitsConstPositiveSecMetric},
  both built on the $\Ric$ and $\Sect$ definitions above. Each quantifies
  existentially over a $C^1$ Levi-Civita connection of the metric and tests
  against $C^2$ vector fields; by Theorem \ref{thm:ricci-well-defined} the
  choice of connection is immaterial, so both are statements about the metric.
  (An earlier version tested sectional curvature against arbitrary fields,
  where $R(X,Y)Y$ is junk, and required no regularity of the connection, where
  $\Ric$ is junk; both were corrected in September 2026.) With
  Definition \ref{def:ricci-of-metric} both hypotheses become statements
  about $\Ric(\g)$ and $\Sect(\g)$ directly:
  \texttt{admitsPositiveRicciMetric\_iff} and
  \texttt{admitsConstPositiveSecMetric\_iff}. It carries no
  \texttt{\textbackslash lean} marker here because \texttt{proof\_wanted}
  produces a private declaration that \texttt{checkdecls} cannot resolve.

  The proof is years away. Stating it was not possible before
  2026-08-13: neither $\Ric$ nor $\Sect$ existed in any Lean library.
\end{theorem}

\chapter{Comparison geometry and non-negative curvature}
\label{chap:comparison}

Morgan--Tian devote their Chapter 2 to this material and it is used throughout
everything below. An earlier version of this blueprint omitted it entirely,
which understated the foundations by a full chapter.

\begin{theorem}[Cheeger--Gromoll splitting theorem]
  \label{thm:splitting}
  A complete manifold of non-negative Ricci curvature containing a line splits
  isometrically as a product with $\R$. \emph{(Morgan--Tian 2.5.)}

  Consumed repeatedly in the classification of ancient solutions, where the
  limits that arise at infinity are shown to split.
\end{theorem}

\begin{theorem}[Soul theorem]
  \label{thm:soul}
  \uses{thm:splitting}
  A complete non-compact manifold of non-negative sectional curvature is
  diffeomorphic to the normal bundle of a compact totally geodesic submanifold.
  \emph{(Morgan--Tian 2.3, via Busemann functions, 2.1--2.2.)}
\end{theorem}

\begin{definition}[$\epsilon$-neck]
  \label{def:neck}
  \uses{def:sectional}
  A region almost isometric, after rescaling, to a cylinder
  $S^2 \times (-\epsilon^{-1}, \epsilon^{-1})$. \emph{(Morgan--Tian 2.6.)}

  The basic local model for where surgery happens. Purely a definition, but a
  fiddly one --- the whole surgery construction is organized around which
  regions are necks and how deep.
\end{definition}

\chapter{Perelman's monotone quantities}
\label{chap:entropy}

Hamilton's hypothesis, positive Ricci curvature, is exactly what makes his
argument work and exactly what a general $3$-manifold does not satisfy. Without
it the flow forms singularities that must be understood rather than avoided,
and understanding them requires ruling out one pathology above all: the metric
collapsing along a sequence of rescalings, so that no limiting model exists.
Perelman's monotone functionals are what rule it out.

\begin{definition}[The $\mathcal{F}$-functional]
  \label{def:f-functional}
  \uses{def:ricci}
  $\mathcal{F}(\g,f) = \int_\M (\scal + |\nab f|^2)\, e^{-f}\, d\vol_\g$.
\end{definition}

\begin{theorem}[Ricci flow is a gradient flow]
  \label{thm:f-monotone}
  \uses{def:f-functional, def:ricci-flow}
  Coupled to the backward heat equation for $f$, the flow is the gradient flow
  of $\mathcal{F}$, and $\mathcal{F}$ is nondecreasing with
  $\partial_t \mathcal{F} = 2\int |\Ric + \nab^2 f|^2 e^{-f}$.

  This retired a long-standing belief that Ricci flow admitted no Lyapunov
  functional, and it is the conceptual heart of Perelman's first preprint.
\end{theorem}

\begin{definition}[The $\mathcal{W}$-entropy]
  \label{def:w-entropy}
  \uses{def:f-functional}
  A scale-invariant modification of $\mathcal{F}$, incorporating a parameter
  $\tau$ with $\dot\tau = -1$ and normalized by
  $\int (4\pi\tau)^{-n/2} e^{-f} d\vol = 1$.
\end{definition}

\begin{theorem}[Monotonicity of $\mathcal{W}$]
  \label{thm:w-monotone}
  \uses{def:w-entropy, thm:f-monotone}
  $\mathcal{W}$ is nondecreasing along the coupled flow.
\end{theorem}

\begin{definition}[$\mathcal{L}$-length, $\mathcal{L}$-geodesics, reduced volume]
  \label{def:reduced-volume}
  \uses{def:ricci-flow}
  The $\mathcal{L}$-length of a path in spacetime, the $\mathcal{L}$-exponential
  map, the reduced length $\ell$, and the reduced volume
  $\tilde V(\tau) = \int \tau^{-n/2} e^{-\ell} d\vol$.
  \emph{(Morgan--Tian Chapter 6, seven sections; Chapter 7 extends it to
  complete flows of bounded curvature.)}

  This is a whole comparison geometry built on the flow rather than on a fixed
  metric --- $\mathcal{L}$-geodesics, an injectivity domain, second-order
  differential inequalities, Lipschitz estimates for $\ell$. An earlier version
  of this blueprint compressed it to a single definition. That was wrong by two
  chapters.
\end{definition}

\begin{theorem}[Monotonicity of the reduced volume]
  \label{thm:reduced-volume-monotone}
  \uses{def:reduced-volume}
  $\tilde V$ is nonincreasing in $\tau$, with equality only on gradient
  shrinking solitons.

  A second, independent monotone quantity. It is the one that does the work in
  the classification of ancient solutions, where $\mathcal{W}$ is not sharp
  enough.
\end{theorem}

\begin{theorem}[No local collapsing]
  \label{thm:no-local-collapsing}
  \uses{thm:w-monotone, def:generalized-flow}
  A flow on a closed manifold on a finite time interval is $\kappa$-noncollapsed
  at every scale below a fixed one: wherever $|\Rm| \le r^{-2}$ on a ball of
  radius $r$, that ball has volume at least $\kappa r^n$.
  \emph{(Morgan--Tian Chapter 8, proved first for generalized flows, then
  specialized to compact ones.)}

  This is what makes rescaling limits exist. Every compactness argument below
  consumes it.
\end{theorem}

\chapter{Singularity models}
\label{chap:kappa}
Blow-up limits at a singularity are the objects this chapter describes, and
before any of them can be recognised one has to know they have non-negative
curvature. In dimension three that is not automatic --- it is the
Hamilton--Ivey estimate, and it was missing from this chart entirely.

\begin{theorem}[Hamilton--Ivey pinching]
  \label{thm:hamilton-ivey}
  \uses{thm:max-tensor, lem:evolution-rm, lem:ivey-ode}
  Let $(\M^3, \g(t))$ be a Ricci flow on a closed three-manifold, normalised so
  that the smallest eigenvalue $\nu$ of the curvature operator satisfies
  $\nu \ge -1$ at $t = 0$. Then at every point and every $t \ge 0$ at which
  $\nu < 0$,
  \[
    \scal \;\ge\; |\nu| \bigl( \log |\nu| + \log(1+t) - 3 \bigr).
  \]
  \emph{(Hamilton 1999 §24; Ivey 1993; Morgan--Tian 4.4; Chow--Knopf 6.44.)}

  Negative curvature is dominated by the scalar curvature, at a rate that
  degenerates only logarithmically. Rescaling a singularity multiplies $\scal$
  by a factor tending to infinity while the estimate is scale-invariant up to
  the logarithm, so every blow-up limit has $\nu \ge 0$: the limits are
  \emph{non-negatively} curved. That is exactly the standing hypothesis in
  Definition~\ref{def:kappa-solution}, and without this theorem no blow-up
  limit is ever known to satisfy it.

  \textbf{Structurally this is Lemma~\ref{lem:pinching-ode} again} --- a closed
  convex set of curvature operators preserved by the ODE
  $\dot{\Rm} = \Rm^2 + \Rm^{\#}$, transported to the flow by the tensor
  maximum principle. Hamilton's set is
  \[
    K : \quad \lambda+\mu+\nu \ge -3
    \quad\text{and}\quad
    \nu + f^{-1}(\lambda+\mu+\nu) \ge 0,
    \qquad f(x) = x(\log x - 3) \ \text{on}\ [e^2,\infty),
  \]
  and the whole content at the ODE level is that $K$ is preserved.
  \emph{(Cao--Zhu, Theorem 2.4.1.)} \textbf{The ODE half is done} ---
  Lemma~\ref{lem:ivey-ode} --- \textbf{and so is everything
  Theorem~\ref{thm:max-tensor} asks of $K$}: it is closed, convex, and
  invariant under the ODE (Lemma~\ref{lem:ivey-convex}). An earlier version of
  this node said the transport still needed $f^{-1}$ on $[-e^2,\infty)$ and its
  concavity; it does not, and never did. What remains is the flow itself,
  $\partial_t \Rm = \Delta \Rm + Q$ (Lemma~\ref{lem:evolution-rm}). For
  Hamilton's later improvement with the $\log(1+t)$ term one additionally needs
  a form of Theorem~\ref{thm:max-tensor} for a \emph{time-dependent} family
  $\{K_t\}$; that is a separate, still-open item.
\end{theorem}

\begin{lemma}[The scalar curvature is nondecreasing along the curvature ODE]
  \label{lem:ivey-scal}
  \lean{RicciFlowBlueprint.Pinching.IsCurvatureODE.le_scal}
  \leanok
  \uses{lem:pinching-ode}
  Along Hamilton's curvature ODE,
  \[
    \tfrac{d}{dt}(\lambda+\mu+\nu)
    = \lambda^2+\mu^2+\nu^2+\lambda\mu+\lambda\nu+\mu\nu
    = \tfrac12\bigl[(\lambda+\mu)^2+(\lambda+\nu)^2+(\mu+\nu)^2\bigr] \ge 0,
  \]
  so every lower bound $\lambda+\mu+\nu \ge c$ is preserved. With $c = -3$ this
  is the first of the two inequalities cutting out $K$.
\end{lemma}

\begin{lemma}[The Hamilton--Ivey boundary inequality]
  \label{lem:ivey-boundary}
  \lean{RicciFlowBlueprint.Pinching.hamiltonIvey_boundary_of_nonneg,
        RicciFlowBlueprint.Pinching.hamiltonIvey_boundary_of_neg}
  \leanok
  \uses{lem:pinching-ode}
  The second inequality cutting out $K$ reads
  $\lambda+\mu \ge (-\nu)[\log(-\nu)-2]$ whenever $\nu \le -e^2$, and need only
  be checked on the boundary. There the defining relation eliminates the
  logarithm and what is left is polynomial. Writing $N = -\nu > 0$ and
  $L = \log N$:
  \begin{itemize}
    \item if $\mu \ge 0$, the boundary relation is $\lambda + \mu = N(L-2)$ and
      the required inequality
      $\dot\lambda + \dot\mu \ge (L-1)\,\dot{\overparen{(-\nu)}}$ becomes,
      after multiplying by $N$,
      \[
        N(\lambda^2+\mu^2) + N^3 + \lambda\mu(\lambda+\mu+N) \ge 0 ;
      \]
    \item if $\mu < 0$, put $P = -\mu$, so the ordering $\nu \le \mu$ is
      $P \le N$ and the boundary relation is $\lambda = P + N(L-2)$; the
      required inequality becomes
      \[
        (\lambda^2 - \lambda P + P^2)(N - P) + P^3 + N^3 \ge 0 .
      \]
  \end{itemize}
  Both hold because every factor is nonnegative --- in the second,
  $\lambda^2-\lambda P+P^2 = (\lambda - P/2)^2 + \tfrac34 P^2$.
  \emph{(Cao--Zhu, Theorem 2.4.1, Cases (i) and (ii); in Lean
  \texttt{hamiltonIvey\_boundary\_of\_nonneg} and
  \texttt{hamiltonIvey\_boundary\_of\_neg}.)}
\end{lemma}

\begin{lemma}[Hamilton's pinching set, and what it encodes]
  \label{lem:ivey-set}
  \lean{RicciFlowBlueprint.Pinching.IsIveyPinched,
        RicciFlowBlueprint.Pinching.isIveyPinched_of_neg_one_le,
        RicciFlowBlueprint.Pinching.le_of_isIveyPinched,
        RicciFlowBlueprint.Pinching.iveyF_le_neg_three}
  \leanok
  \uses{lem:ivey-scal, lem:ivey-boundary}
  Because $f^{-1}$ takes values in $[e^2,\infty)$, the second condition cutting out $K$ holds
  automatically when $-\nu \le e^2$ and is $f(-\nu) \le \lambda+\mu+\nu$ otherwise. So $K$ can
  be written with no inverse function at all:
  \[
    K : \quad \lambda+\mu+\nu \ge -3
    \quad\text{and}\quad
    \bigl(\, -\nu \le e^2 \ \text{ or } \ f(-\nu) \le \lambda+\mu+\nu \,\bigr).
  \]
  In that form both ends of Hamilton's argument are elementary.

  \emph{Entry.} An ordered triple with $\nu \ge -1$ lies in $K$: the ordering gives
  $\lambda+\mu+\nu \ge 3\nu \ge -3$, and $-\nu \le 1 \le e^2$.

  \emph{Exit.} An ordered triple in $K$ with $\nu < 0$ satisfies
  $\scal \ge (-\nu)(\log(-\nu) - 3)$. Write $N = -\nu$. For $N > e^2$ this is the second
  condition verbatim. For $N \le 1$ we have $\log N \le 0$, so $f(N) \le -3N$, and the ordering
  gives $\scal \ge 3\nu = -3N$. For $1 \le N \le e^2$ the function $f$ is antitone --- there
  $f' = \log N - 2 \le 0$ --- so $f(N) \le f(1) = -3$, and the \emph{first} condition of $K$
  finishes it. That middle range is the whole reason $K$ carries the otherwise unmotivated
  bound $\lambda+\mu+\nu \ge -3$.
\end{lemma}

\begin{lemma}[The pinching set is closed and convex]
  \label{lem:ivey-convex}
  \lean{RicciFlowBlueprint.Pinching.iveyG,
        RicciFlowBlueprint.Pinching.convexOn_iveyG,
        RicciFlowBlueprint.Pinching.isIveyPinched_iff_iveyG,
        RicciFlowBlueprint.Pinching.iveyPinchedSet,
        RicciFlowBlueprint.Pinching.convex_iveyPinchedSet,
        RicciFlowBlueprint.Pinching.isClosed_iveyPinchedSet,
        RicciFlowBlueprint.Pinching.IsCurvatureODE.isIveyPinched,
        RicciFlowBlueprint.Pinching.IsCurvatureODE.mem_iveyPinchedSet}
  \leanok
  \uses{lem:ivey-set, thm:max-tensor}
  $K$ is a closed convex subset of $\R^3$ --- which is what
  Theorem \ref{thm:max-tensor} requires of it, and the one thing missing before
  Lemma \ref{lem:ivey-ode} can be carried from the ODE to the flow.

  \textbf{Proved} (\texttt{IveyConvex.lean}), and \textbf{without building
  $f^{-1}$.} The literature restores convexity by rewriting the second condition
  as $-\nu \le f^{-1}(\lambda+\mu+\nu)$ and invoking concavity of $f^{-1}$; an
  earlier version of these notes recorded that inverse, and its concavity, as
  the remaining obstacle. It is not needed. The disjunction of
  Lemma \ref{lem:ivey-set} collapses to a single inequality against
  \[
    G(\nu) \;=\; f\bigl(\max(-\nu,\, e^2)\bigr),
  \]
  because $f$ is increasing on $[e^2,\infty)$: when $-\nu \le e^2$ we get
  $G(\nu) = f(e^2) = -e^2 \le -3$, which the first condition of $K$ already
  supplies, and otherwise $G(\nu) = f(-\nu)$ on the nose
  (\texttt{isIveyPinched\_iff\_iveyG}).
\end{lemma}

\begin{proof}
  \leanok
  $f$ is convex on $[e^2,\infty)$ because $f' = \log x - 2$ is monotone there,
  and increasing there because $f' \ge 0$; $\nu \mapsto \max(-\nu, e^2)$ is
  convex, being a pointwise maximum of two affine functions, and its image is
  exactly $[e^2,\infty)$. A convex function composed with a convex function that
  lands where the outer one is monotone is convex, so $G$ is convex on all of
  $\R$ (\texttt{convexOn\_iveyG}). The set is then cut out by three half-spaces
  --- $\nu \le \mu$, $\mu \le \lambda$, $\lambda+\mu+\nu \ge -3$ --- and the
  single convex inequality $G(\nu) \le \lambda+\mu+\nu$, hence is convex; and
  closed, $G$ being continuous.

  \emph{Invariance, in the same language.} An ordered solution of the curvature
  ODE normalised by $\nu(0) \ge -1$ stays in $K$ for all time
  (\texttt{IsCurvatureODE.isIveyPinched}, and
  \texttt{IsCurvatureODE.mem\_iveyPinchedSet} as membership of the subset of
  $\R^3$): $R$ is nondecreasing and starts at $\ge 3\nu(0) \ge -3$, and where
  $\nu < 0$ the second condition is Lemma \ref{lem:ivey-ode} verbatim, while
  where $\nu \ge 0$ the branch $-\nu \le e^2$ is free. So both of the inputs
  Theorem \ref{thm:max-tensor} needs about $K$ --- closed convex, and invariant
  under the ODE --- are now available; what is still missing for the transport
  is the flow itself, $\partial_t \Rm = \Delta \Rm + Q$
  (Lemma \ref{lem:evolution-rm}).
\end{proof}

\begin{lemma}[Hamilton--Ivey for the curvature ODE]
  \label{lem:ivey-ode}
  \lean{RicciFlowBlueprint.Pinching.IsCurvatureODE.hamiltonIvey,
        RicciFlowBlueprint.Pinching.IsCurvatureODE.iveyPsi_nonneg,
        RicciFlowBlueprint.Pinching.iveyE_nonneg,
        RicciFlowBlueprint.Pinching.IsCurvatureODE.nonneg_preserved}
  \leanok
  \uses{lem:pinching-ode, lem:ivey-scal, lem:ivey-set, lem:ivey-convex}
  Let $\lambda \ge \mu \ge \nu$ solve Hamilton's curvature ODE on $[0,T]$ with
  $\nu(0) \ge -1$. Then at every time at which $\nu < 0$,
  \[
    \lambda + \mu + \nu \ \ge\ (-\nu)\bigl(\log(-\nu) - 3\bigr).
  \]

  Cao--Zhu check invariance of $K$ on its boundary, in two cases according to the
  sign of $\mu$. The two cases are the same inequality: the hypothesis of the
  second, $\lambda = -\mu + N(L-2)$ with $N = -\nu$, is the hypothesis of the
  first, $\lambda+\mu = N(L-2)$, and eliminating $L$ from either leaves the same
  polynomial
  \[
    \texttt{iveyE} \;=\; N(\lambda^2+\mu^2) + N^3 + \lambda\mu(\lambda+\mu+N),
  \]
  which is nonnegative under the ordering alone --- no boundary relation, no
  logarithm. So the boundary argument becomes a linear Gr\"onwall comparison: for
  the defect $\Psi = \scal - f(-\nu)$ one has the identity
  \[
    \Psi' \;=\; \frac{\texttt{iveyE}}{N} \;-\; \Psi\,\frac{N^2 + \lambda\mu}{N},
  \]
  valid wherever $\nu < 0$, so $\Psi' \ge a\Psi$ with $a = -(N^2+\lambda\mu)/N$
  continuous, and $\Psi \ge 0$ is preserved.

  That $\nu < 0$ on the whole interval is not an extra hypothesis: non-negative
  curvature is preserved by the ODE ($\dot\nu \ge \lambda\nu$ under the ordering),
  so a $\nu$ that is negative at $t$ was negative on all of $[0,t]$, and where
  $\nu \ge 0$ the estimate is vacuous.
\end{lemma}

\begin{theorem}[Geometric convergence of manifolds and of flows]
  \label{thm:compactness}
  \uses{thm:shi, thm:no-local-collapsing}
  A sequence of pointed Ricci flows with uniformly bounded curvature and a
  uniform injectivity-radius lower bound at the basepoints subconverges, in
  $C^\infty$ on compact sets, to a limiting pointed flow. Blow-up limits at a
  singularity are obtained this way. \emph{(Morgan--Tian Chapter 5, five
  sections: convergence of manifolds, of flows, Gromov--Hausdorff, blow-up
  limits, splitting limits at infinity.)}

  \textbf{Blocked, separately from the PDE gap.} Mathlib has the
  Gromov--Hausdorff distance between compact metric spaces; the pointed,
  smooth, curvature-bounded convergence theory for manifolds is a different and
  much larger object, and none of it exists. It is the second independent
  analytic prerequisite on this road and the first not shared with Hamilton's
  theorem.
\end{theorem}

\begin{definition}[$\kappa$-solution]
  \label{def:kappa-solution}
  \uses{thm:no-local-collapsing, thm:hamilton-ivey}
  An ancient solution --- defined on $(-\infty, 0]$ --- that is
  $\kappa$-noncollapsed, has bounded non-negative curvature operator, and is not
  flat. These are exactly the models that arise from rescaling a singularity.
\end{definition}

\begin{theorem}[Classification of three-dimensional $\kappa$-solutions]
  \label{thm:kappa-classification}
  \uses{def:kappa-solution, thm:compactness, thm:reduced-volume-monotone,
        thm:splitting}
  In dimension three every $\kappa$-solution is, at every point and scale, close
  to one of a short explicit list: the round shrinking $S^3/\Gamma$, the round
  shrinking neck $S^2 \times \R$, or a capped variant.
  \emph{(Morgan--Tian Chapter 9, eight sections --- the asymptotic gradient
  shrinking soliton, splitting at infinity, classification of gradient shrinking
  solitons in dimensions 2 and 3, a universal $\kappa$, asymptotic volume,
  compactness of the space of $3$-dimensional $\kappa$-solutions.)}

  The technical core of Perelman's second preprint and the longest single
  chapter of the book.
\end{theorem}

\begin{theorem}[Bounded curvature at bounded distance]
  \label{thm:bounded-curvature}
  \uses{thm:kappa-classification, thm:compactness}
  Curvature control at a point propagates to control on a definite
  neighbourhood. \emph{(Morgan--Tian Chapter 10.)}

  A separate chapter, proved by contradiction through an incomplete geometric
  limit and a comparison of Gromov--Hausdorff and smooth limits. Omitted
  entirely from an earlier version of this chart.
\end{theorem}

\begin{theorem}[Canonical neighborhood theorem]
  \label{thm:canonical-nbhd}
  \uses{thm:kappa-classification, thm:bounded-curvature, def:neck}
  In a flow on a closed $3$-manifold, every point of sufficiently large
  curvature has a neighborhood that, after rescaling, is close to a
  corresponding piece of a $\kappa$-solution --- a neck, a cap, or a closed
  spherical piece. \emph{(Morgan--Tian Chapter 11 supplies the geometric limits
  of generalized flows this rests on; the appendix, Chapter 19, supplies the
  neck and cap geometry.)}

  This converts an analytic classification into a topological description of
  where and how the manifold is about to pinch, and it is what makes surgery
  well-defined rather than arbitrary.
\end{theorem}

\chapter{Ricci flow with surgery}
\label{chap:surgery}

Here the object being studied changes character: the manifold itself is cut and
reglued as the flow proceeds. The encouraging discovery on reading Morgan--Tian's
structure is that they do \emph{not} treat this as an informal procedure. The
flow is redefined on a single geometric object --- a surgery space-time --- and
the surgery process becomes an equation on that object. That is precisely the
kind of reformulation formalization needs, and it moderates an earlier claim
here that surgery is unformalizable in principle. It remains unattempted
anywhere, and it remains long.

\begin{definition}[The standard solution]
  \label{def:standard-solution}
  \uses{def:kappa-solution}
  The explicit flow on $\R^3$ with rotationally symmetric initial metric,
  asymptotic to a cylinder, used as the model for the cap glued in at surgery.
  \emph{(Morgan--Tian Chapter 12, seven sections.)}

  \textbf{Carries its own PDE prerequisite.} Uniqueness of the standard solution
  is proved through the \emph{harmonic map flow} (12.5), a second geometric flow
  with its own existence theory. This road needs more than one parabolic theory.
\end{definition}

\begin{theorem}[Surgery on a $\delta$-neck]
  \label{thm:neck-surgery}
  \uses{def:neck, def:standard-solution}
  The metric on a sufficiently deep neck may be replaced by a capped metric with
  controlled curvature, preserving the pinching hypotheses.
  \emph{(Morgan--Tian Chapter 13.)}
\end{theorem}

\begin{definition}[Surgery space-time and Ricci flow with surgery]
  \label{def:surgery}
  \uses{thm:canonical-nbhd, thm:neck-surgery, def:generalized-flow}
  A space-time whose time-slices are the evolving manifolds, singular along the
  surgery caps, carrying a horizontal metric satisfying the generalized Ricci
  flow equation. \emph{(Morgan--Tian Chapter 14.)}
\end{definition}

\begin{theorem}[Existence of Ricci flow with surgery]
  \label{thm:surgery-exists}
  \uses{def:surgery, thm:canonical-nbhd, thm:no-local-collapsing}
  Let $(\M, \g_0)$ be a closed Riemannian $3$-manifold containing no embedded,
  locally separating $\R P^2$. Then there is a Ricci flow with surgery defined
  for all $t \in [0,\infty)$ with initial metric $\g_0$, whose discontinuity
  times form a discrete subset of $[0,\infty)$, and where the topological change
  across a surgery time is a connected sum decomposition together with removal
  of components diffeomorphic to $S^2 \times S^1$, $\R P^3 \# \R P^3$, the
  non-orientable $S^2$-bundle over $S^1$, or a manifold of constant positive
  curvature. \emph{(Morgan--Tian Theorem 0.3; Chapters 15--17.)}

  The $\R P^2$ hypothesis is not decoration and was missing from an earlier
  version of this chart. The proof is a mutual induction over surgery times:
  parameters are chosen in advance, and non-collapsing (Chapter 16) and the
  canonical neighborhood assumption (Chapter 17) are re-established after every
  surgery. Chapters 15--17 are the bulk of the book.
\end{theorem}

\chapter{The spherical space form theorem}
\label{chap:spaceform}

This is the terminal node. The route reaches it \emph{without} proving the
geometrization conjecture, which matters a great deal for how much has to be
formalized.

\begin{theorem}[Finite-time extinction]
  \label{thm:finite-extinction}
  \uses{thm:surgery-exists}
  Let $\M$ be a closed $3$-manifold whose fundamental group is a free product of
  finite groups and infinite cyclic groups. Then $\M$ admits no locally
  separating $\R P^2$, so a Ricci flow with surgery exists for all positive time,
  and it becomes extinct after some $T < \infty$: the time-slices $\M_t$ are
  empty for $t \ge T$. \emph{(Morgan--Tian Theorem 0.4, proved in Chapter 18.)}

  \textbf{Carries a third PDE prerequisite.} Morgan--Tian follow Perelman's
  third preprint and route the proof through the \emph{curve-shrinking flow}
  (18.4), yet another geometric flow with its own existence and regularity
  theory. Chapter 18 is seven sections and largely self-contained analysis.
\end{theorem}

\begin{theorem}[Decomposition]
  \label{thm:decomposition}
  \uses{thm:surgery-exists, thm:finite-extinction}
  A closed $3$-manifold whose fundamental group is a free product of finite
  groups and infinite cyclic groups is a connected sum of spherical space forms,
  copies of $S^2 \times S^1$, and copies of the non-orientable $S^2$-bundle over
  $S^1$. \emph{(Morgan--Tian Theorem 0.1, deduced from 0.3 and 0.4 by downward
  induction on the surgery times.)}
\end{theorem}

\begin{theorem}[Spherical space form / elliptization]
  \label{thm:space-form}
  \uses{thm:decomposition, def:sectional}
  A closed $3$-manifold with finite fundamental group admits a metric of
  constant positive sectional curvature, and is therefore diffeomorphic to
  $S^3/\Gamma$ for a finite $\Gamma \subset SO(4)$ acting freely.

  A finite group is in particular a free product of finite groups, so
  Theorem~\ref{thm:decomposition} applies. An $S^2 \times S^1$ summand would
  contribute a $\Z$ factor and a connected sum of two non-trivial groups is
  infinite, so with $\pi_1$ finite exactly one spherical summand survives.
\end{theorem}

\begin{corollary}[Poincar\'e conjecture]
  \label{cor:poincare}
  \uses{thm:space-form}
  A simply connected closed $3$-manifold is diffeomorphic to $S^3$. The case
  $\Gamma = 1$.

  Mathlib already states this, unproved, in
  \texttt{Mathlib/Geometry/Manifold/PoincareConjecture.lean}. It is the only
  node on this road that Lean has so far looked at.
\end{corollary}

\section*{What this route deliberately avoids}

Geometrization requires, in addition to everything above, the analysis of the
\emph{long-time} behaviour of the flow: the thick--thin decomposition,
collapsing with local lower curvature bounds, and the recognition of graph
manifolds --- together with the JSJ decomposition and Thurston's eight
geometries on the topological side. None of that is needed here. Finite
extinction short-circuits it: the flow simply ends, so its long-time behaviour
never has to be described. Morgan--Tian say so explicitly --- the argument
``avoids all references to the nature of the flow as time goes to infinity'' ---
and they defer geometrization to a second volume published seven years later.

\section*{Calibration against Morgan--Tian}

Morgan--Tian is 19 chapters and roughly 100 sections over 492 pages, all of it
downstream of the material in Chapters \ref{chap:curvature} and
\ref{chap:levi-civita} here, none of which they prove --- their Chapter 1 is
six sections of Riemannian geometry assumed as background, and it is the part
Mathlib does not have.

Reading their structure corrected five errors in this chart: an entire chapter
of comparison geometry was missing (Chapter \ref{chap:comparison}), the
$\mathcal{L}$-geometry was compressed from two chapters into one definition,
bounded-curvature-at-bounded-distance was absent, the $\R P^2$ hypothesis
was dropped from the surgery theorem, and the Hamilton--Ivey estimate
(Theorem \ref{thm:hamilton-ivey}) --- without which no blow-up limit is known
to be non-negatively curved, and so no $\kappa$-solution is ever produced ---
did not appear at all. It also produced one finding that no
citation list would have surfaced: this road needs \emph{three} distinct
parabolic theories --- Ricci flow, harmonic map flow, and curve-shrinking flow
--- not one.
